%% qcert.tex -- QCert paper draft
%% "QCert: Verifying Sixteen Quantum Circuit Completeness Results
%% in One Agda Framework"
%%
%% Shares styles, figures, and bibliography with ../planqc/.
%%
%% Build from pldi/: latexmk -pdf qcert.tex
%% `anonymous' for double-blind review; remove it for the camera-ready version.
% \documentclass[acmsmall,nonacm,11pt,screen]{acmart}
\documentclass[acmsmall,nonacm,11pt,screen,anonymous]{acmart}

\settopmatter{printacmref=false, printccs=false, printfolios=true}
%% No ``Anon.'' in the running head of even pages (double-blind review);
%% remove this with the `anonymous' option for the camera-ready version.
\makeatletter
\def\@shortauthors{}
\makeatother
\setcopyright{none}
\renewcommand\footnotetextcopyrightpermission[1]{}

%% The contact footnote lists the e-mail addresses only; the names and
%% affiliations stay in the title block.
%% Commented out for double-blind review (it prints ``Contact:'' even
%% with no authors); restore with the authors for the camera-ready version.
% \makeatletter
% \def\@mkauthorsaddresses{%
%   Contact:
%   \bgroup
%   \def\position##1{\unskip\ignorespaces}%
%   \def\institution##1{\unskip\ignorespaces}%
%   \def\city##1{\unskip\ignorespaces}%
%   \def\state##1{\unskip\ignorespaces}%
%   \renewcommand\department[2][0]{\unskip\ignorespaces}%
%   \def\country##1{\unskip\ignorespaces}%
%   \def\and{\unskip, }%
%   \def\@author##1{\ignorespaces}%
%   \def\email##1##2{\unskip\ \mbox{\nolinkurl{##2}}}%
%   \addresses
%   \egroup}
% %% No final period, which would read as part of the last address.
% \def\@setauthorsaddresses{\@authorsaddresses\unskip}
% \makeatother

\usepackage{mathtools}
\usepackage{booktabs}
\usepackage{tabularx}
\newcolumntype{Y}{>{\raggedright\arraybackslash}X}
%% This file lives in pldi/; the style file, the figures
%% and the bibliography are in ../planqc/.
\makeatletter
\def\input@path{{../planqc/}}
\makeatother
\usepackage{agda-style}
%% Agda's own style file, for the code blocks that `agda --latex'
%% generated for the qudit Clifford formalization (qudit/agda/).
\usepackage{agda}
\renewcommand{\AgdaCodeStyle}{\footnotesize}
\setlength{\mathindent}{1em}
%% Characters of those code blocks that agda-style.sty does not declare.
\DeclareUnicodeCharacter{02E2}{\textsuperscript{s}}
\DeclareUnicodeCharacter{1D47}{\textsuperscript{b}}
\DeclareUnicodeCharacter{1D48}{\textsuperscript{d}}
\DeclareUnicodeCharacter{1D49}{\textsuperscript{e}}
\DeclareUnicodeCharacter{1D56}{\textsuperscript{p}}
\DeclareUnicodeCharacter{1D5B}{\textsuperscript{v}}
\DeclareUnicodeCharacter{1DA0}{\textsuperscript{f}}
\DeclareUnicodeCharacter{1D62}{\textsubscript{i}}
\DeclareUnicodeCharacter{208B}{\textsubscript{-}}
\DeclareUnicodeCharacter{2190}{\ensuremath{\leftarrow}}
\DeclareUnicodeCharacter{2193}{\ensuremath{\downarrow}}
\DeclareUnicodeCharacter{21E3}{\ensuremath{\downarrow}}
\DeclareUnicodeCharacter{224B}{\ensuremath{\approx\!\!\approx}}
\DeclareUnicodeCharacter{00B1}{\ensuremath{\pm}}
\usepackage{subcaption}
\usepackage{cleveref}

%% TikZiT figures. A subset of tikzit.sty (the layer stack, the blank
%% node style and \tikzfig), which would load enumitem, which clashes
%% with acmart. proofs.tikzstyles has the styles of ../planqc/figures/*.tikz.
\usepackage{tikz}
\usetikzlibrary{backgrounds,arrows,shapes,shapes.geometric,shapes.misc}
\pgfdeclarelayer{edgelayer}
\pgfdeclarelayer{nodelayer}
\pgfsetlayers{background,edgelayer,nodelayer,main}
\tikzset{tikzit fill/.initial=0,tikzit draw/.initial=0,tikzit shape/.initial=0,tikzit category/.initial=0}
\tikzstyle{none}=[inner sep=0mm]
\tikzstyle{tikzfig}=[baseline=-0.25em,scale=0.7]
\newcommand{\tikzfig}[1]{{\tikzstyle{every picture}=[tikzfig]\input{#1.tikz}}}
\input{../planqc/proofs.tikzstyles}
%% The circuits of the qudit Clifford formalization (qudit/figures/),
%% with their own styles, scaled by a fixed factor so that gate boxes
%% stay proportionate across a grid of circuits.
\usetikzlibrary{arrows.meta}
\definecolor{normGreen}{RGB}{214, 243, 221}
\definecolor{normRed}{RGB}{249, 196, 196}
\definecolor{normYellow}{RGB}{253,246,208}
\definecolor{zxgreen}{RGB}{216,248,216}
\definecolor{zxred}{RGB}{221,165,165}
\definecolor{zxhad}{RGB}{255,255,0}
\definecolor{zxpaulibox}{RGB}{138,159,198}
\input{qudit/circuits.tikzstyles}
\tikzstyle{qtikzfig}=[baseline=-0.25em,scale=0.38]
\newcommand{\ufig}[2][1]{$\vcenter{\hbox{\scalebox{#1}{{\tikzstyle{every picture}=[qtikzfig]\input{qudit/figures/#2.tikz}}}}}$}
%% The comparison of the Agda and Lean formalizations (figures/compare.tex,
%% generated by compare/mkfig.py). Agda blue and Lean orange are the first
%% two slots of a colour-blind-safe categorical palette; Agda's bar is above
%% Lean's and the markers differ, so the figure also reads in greyscale.
\usepackage{pgfplots}
\pgfplotsset{compat=1.18}
\usepgfplotslibrary{groupplots}
\usetikzlibrary{patterns,calc}
\definecolor{cmpAgda}{HTML}{2A78D6}
\definecolor{cmpLean}{HTML}{EB6834}
\colorlet{cmpAxis}{black!55}
\colorlet{cmpGrid}{black!12}
\tikzset{cmpLink/.style={black!30, line width=0.7pt}}

\citestyle{acmnumeric}
%% Table 1 keeps author-year citations: ``Greylyn [2014]''.
\newcommand{\tabcite}[1]{\citeauthor{#1}~[\citeyear{#1}]}
%% A remark, in the style of acmart's definitions (defined after acmart's
%% own theorem environments, so that it can share their counter).
\AtEndPreamble{\theoremstyle{acmdefinition}\newtheorem{remark}[theorem]{Remark}}
%% Named ``Bian, Li et al.'' (authors in alphabetical order), as a link.
\defcitealias{bian2026qupit}{Bian, Li et al.}
\newcommand{\quditcite}{\citetalias{bian2026qupit}~\cite{bian2026qupit}}
\newcommand{\rewrites}[1]{\overset{#1^{*}}{\longleftrightarrow}}

%% Inline Agda identifiers, in a sans-serif face.
\newcommand{\aid}[1]{\mbox{\textsf{\small #1}}}
%% A numbered, referenceable line of quoted Agda: end the line of an
%% agdacode block with escapechar=@ in @\agdatag{label}@.
\newcounter{agdastmt}
\crefname{agdastmt}{statement}{statements}
\Crefname{agdastmt}{Statement}{Statements}
\creflabelformat{agdastmt}{(#2#1#3)}
\newcommand{\agdatag}[1]{\hfill\normalfont\normalsize
  \refstepcounter{agdastmt}(\theagdastmt)\label{#1}}
\newcommand{\Z}{\mathbb{Z}}
\newcommand{\D}{\mathbb{D}}
%% Notation of the qudit Clifford formalization (\Cref{sec:qudit}).
\newcommand{\Zp}{\ensuremath{\mathbb{Z}_p}}
\newcommand{\Sp}[1][n]{\ensuremath{\mathrm{Sp}(2#1,\mathbb{Z}_p)}}
\newcommand{\Pauli}[1][n]{\ensuremath{\mathcal{P}_{#1}}}
\newcommand{\Cliff}[1][n]{\ensuremath{\mathcal{C}_{#1}}}
\newcommand{\PCliff}[1][n]{\ensuremath{\overline{\mathcal{C}}_{#1}}}
\newcommand{\PPauli}[1][n]{\ensuremath{\overline{\mathcal{P}}_{#1}}}
\newcommand{\ket}[1]{\ensuremath{|#1\rangle}}
\newcommand{\sem}[1]{\ensuremath{[\![#1]\!]}}
\newcommand{\ML}{\ensuremath{\mathit{ML}}}

\begin{document}

\title[QCert: Verifying and Extending Quantum Circuit Completeness in Agda and Lean]{QCert: Verifying and Extending Quantum Circuit Completeness in Agda and Lean}
%QCert: Verifying Sixteen Quantum Circuit Completeness Results in One Agda Framework

%% Authors omitted for double-blind review; restore for the camera-ready version.
% \author{Xiaoning Bian}
% \affiliation{%
%   \institution{Tsinghua University}
%   \city{Beijing}
%   \country{China}}
% \email{bian@dal.ca}
% \author{Yuan Feng}
% \affiliation{%
%   \institution{Tsinghua University}
%   \city{Beijing}
%   \country{China}}
% \email{yuan\_feng@tsinghua.edu.cn}
% \author{Sarah Meng Li}
% \affiliation{%
%   \institution{University of Amsterdam and QuSoft}
%   \city{Amsterdam}
%   \country{The Netherlands}}
% \email{sarah.li@uva.nl}
% \author{John van de Wetering}
% \affiliation{%
%   \institution{University of Amsterdam and QuSoft}
%   \city{Amsterdam}
%   \country{The Netherlands}}
% \email{john@vdwetering.name}

\begin{abstract}
% An equational theory for a family of quantum circuits is a set of rewrite rules. The theory 

A set of circuit equations is complete if every equation between circuits denoting the same linear map follows from it, and minimal if no equation can be derived from the others. Establishing these properties takes long case analyses spanning hundreds of pages, making the details hard to check by hand. To address this, we extend an existing formalization library to QCert in Agda and Lean. In QCert, humans set up a unified framework of circuits, matrix semantics, and group presentations, and state each paper's main theorem in it. LLM coding agents then write the proofs following the papers' arguments, and the proof assistant checks them. A human therefore need only review each paper's relations, gate semantics, and theorem statement. Using QCert, we first check sixteen published results, fifteen on completeness and one on minimality. They cover discrete gate sets such as Clifford, Clifford+T, and Toffoli--Hadamard, on qubits or qudits of every odd prime dimension. We verify fourteen of them directly and correct an error in each of the other two. We then prove a new result, generalizing the completeness of Toffoli--Hadamard circuits from three to any number of qubits. We expect QCert to help with other open completeness problems, such as those for quantum circuits in higher dimensions.
%which include the main circuit completeness theorems to date.
\end{abstract}

\maketitle


\section{Introduction}\label{sec:intro}
Rewriting of quantum circuits is a ubiquitous tool in circuit optimization~\cite{hietala2021voqc,nam2018automated,maslov2008simplification,xu2022quartz}, equivalence checking~\cite{yamashita2010fast}, hardware mapping~\cite{sivarajah2020tket}, and reasoning about quantum programs~\cite{fang2026hadamard}. This relies on identifying a set of rewrite rules for a family of quantum circuits that transform a circuit into an equivalent one.
Such a theory should be \emph{sound} so that every rewrite preserves the underlying linear map. It should also be \emph{complete} so that any equivalence between circuits can be derived. Additionally, we want it to be \emph{minimal}, with no redundant rules. 
% Complete theories also support reasoning about quantum programs. For example,
% \mbox{Hadamard-$\Pi$}~\cite{fang2026hadamard} is a complete
% equational theory for reversible quantum programs. 

Completeness results have been established for a wide range of circuit families, including reversible linear circuits~\cite{lafont2003boolean,sakashita2013cnot}, reversible circuits over multiple-control Toffoli gates~\cite{feng2026complete,feng2026ancillas}, Clifford~\cite{selinger2015clifford,makary2021real,li2025qutrit,bian2026qupit}, CNOT-dihedral~\cite{amy2018cnotdihedral,blake2026phaseaffine}, two-qubit Clifford+T~\cite{bian2023cliffordt2}, three-qubit Clifford+CS~\cite{bian2023cliffordcs3}, Toffoli--Hadamard~\cite{amy2024toffoli}, and real Clifford+CH~\cite{clement2026realcliffordch} circuits, as well as universal quantum circuits~\cite{clement2023complete,clement2024extensions}. Completeness has also been obtained for the group of two-level matrices underlying several of these circuit families~\cite{greylyn2014,bian2021un,li2021on,fang2026hadamard,clement2019thesis}. Minimal theories have been given for six of these circuit fragments~\cite{blake2026simpler} and for all quantum circuits~\cite{clement2024minimal}.

Completeness and minimality proofs typically require long case distinctions over gates, relations, and normal forms, whose supporting derivations sometimes span hundreds of pages~\cite{fang2026hadamard,li2025qutrit,clement2026realcliffordch}. In particular, relation reductions require intricate circuit manipulations and repeated rewrites, making them prone to human mistakes. For example, circular reasoning can lead us to remove a rule whose derivation still relies on that rule, hidden in a chain of other derivations. Errors can then propagate through dependent derivations, invalidating the entire reduction and making the error source hard to trace by hand. 

The length of these proofs and their error-prone nature make it desirable to formally verify their correctness using a proof assistant like Lean or Agda. This would not have been feasible in the past due to  the sheer number of lemmas and lengthy rewrites in these papers. Recent LLMs can write correct code for proof assistants, which brings this closer to reality. For example, Anthropic reports that LLMs wrote a 13-million-line Lean proof of Fermat's Last Theorem within 11 days~\cite{anthropic2026flt}. However, care must still be taken. It is important to manually verify the \emph{alignment} between the machine-verified statement and the intended natural-language result. For instance, \cite{bastounis2026navier} showed that a claimed verification of OpenAI's natural-language proof about the Navier--Stokes equations contained mistranslations, so the formal verification was verifying a weaker result by a different proof. A way to solve this is to use a hybrid approach, where humans write and carefully review the definitions and theorem statements~\cite{hales2026reliability}, coding agents write the proofs, and Agda checks them.

Related human-supervised, AI-assisted workflows have been used to formalize quantum program logics~\cite{ren2026qedia}, algorithm correctness and resource bounds~\cite{zhang2026leanquantumalg}, and the quantum soundness of a low individual-degree test underlying $\mathrm{MIP}^{*}=\mathrm{RE}$~\cite{lu2026formalflow}. Lean-QuantumAlg also makes explicit the human review of problem specifications and execution models, while FormalFlow corrects side conditions and intermediate error bounds in the published proof.

Beyond these workflows, formal verification of quantum computing has rarely addressed completeness. Verified compilers and program logics~\cite{paykin2017qwire,hietala2021voqc,lehmann2025vyzx,zhou2023coqq,chareton2021qbricks,li2022vqo,peduri2025qbc} establish program correctness or the soundness of circuit transformations, rather than the completeness of circuit equations. The few mechanized completeness results, for $\Pi$ as a free symmetric rig groupoid~\cite{choudhury2022symmetries} and its square-root extension~\cite{carette2024squareroots}, are coherence arguments. Other recent formalizations address quantum information theory and cryptographic security. Lean-Quantum~\cite{kasaura2026leanquantum} provides operator-theoretic infrastructure and formalizes data processing for the sandwiched R\'enyi relative entropy, and Lean-QIT~\cite{zhu2026leanqit} formalizes quantum source coding and classical communication capacity theorems. FVQKD~\cite{deng2026fvqkd} mechanizes composable finite-key security of BB84 in Rocq, and \citet{benshahar2026cryptography} give an LLM-assisted Lean formalization of composable finite-key security for entanglement-based BB84. These developments share QCert's concern with reusable foundations and faithful specifications, while QCert targets the completeness and minimality of circuit equational theories. General libraries of formalized mathematics, such as Mathlib~\cite{mathlib2020,mathlib-doc-presented}, the AFP~\cite{breitner2010freegroups,kharim2023freegroups}, and MathComp~\cite{gonthier2013oddorder}, also have no verified coset enumeration against a given concrete group, which several completeness proofs in QCert rely on.

\paragraph{Our contributions} In this work we present QCert, a framework in Agda~\cite{norell2007thesis,bove2009agda} and Lean~\cite{moura2021lean4} for representing and verifying quantum circuit completeness and minimality results, and use it to verify the correctness of sixteen completeness results.
% This is based on our formalization of qudit Clifford completeness~\cite{bian2026qupit}, where we built the framework~\cite{bian2027popl} by hand and carried out the main verification steps by hand.
This is based on our formalization of the qudit Clifford completeness result of~\quditcite{} (\Cref{sec:qudit}), where the framework~\cite{bian2027popl} was built and the main verification steps were carried out by hand.
Since all sixteen completeness results use the same framework of quantum circuits, gates, and matrices over a finitely generated ring, this makes it possible to efficiently check the alignment of the formalized results with the natural-language statements from the original papers (\Cref{fig:qcert-pipeline}). Of these sixteen results, we verify fourteen directly and correct an error in each of the other two (\Cref{sec:findings}).

% Draft assumption: this paragraph describes the intended completed state.
% New-result development is ongoing; formal verification remains pending;
% confirm the claims below once both are complete.
QCert also provides a foundation for LLM-assisted exploration of new completeness results. Its shared definitions and verified libraries allow coding agents to build on existing mathematics and automate derivations that would otherwise require substantial manual effort. As a case study, we use this workflow to extend Toffoli--Hadamard completeness from three to any number of qubits (\Cref{sec:toffoli-hadamard}). Alongside the proofs checked in QCert, we are developing a clear mathematical exposition of the underlying arguments. This case study demonstrates QCert's potential as a research tool: it makes automated proof development practical while grounding confidence in machine checking and human review of the definitions and theorem statements.

\begin{figure}[t]
\centering
% Section 1 overview. All artwork is native TikZ and remains editable.
\begingroup
\definecolor{pipelineBlue}{HTML}{205886}
% Cool emerald and mint for human validation; light lemon for exploration.
\definecolor{pipelineGreen}{HTML}{32866B}
\definecolor{pipelineGreenText}{HTML}{23674F}
\definecolor{pipelineMint}{HTML}{EFF9F3}
\definecolor{pipelineMintAccent}{HTML}{BCE5CE}
\definecolor{pipelineLemon}{HTML}{FFF8D1}
\definecolor{pipelineYellow}{HTML}{D6BA2B}
\definecolor{pipelineYellowAccent}{HTML}{F3D34A}
\definecolor{pipelineInk}{HTML}{233447}
\begin{tikzpicture}[
  x=1cm,y=1cm,
  line cap=round,line join=round,
  pipe arrow/.style={-{Stealth[length=1.6mm,width=1.5mm]},
    draw=pipelineInk,line width=.85pt},
  extension arrow/.style={-{Stealth[length=1.5mm,width=1.4mm]},
    draw=pipelineYellow,line width=.8pt,densely dashed},
  alignment arrow/.style={{Stealth[length=1.9mm,width=1.7mm]}-{Stealth[length=1.9mm,width=1.7mm]},
    draw=pipelineGreen,line width=1.3pt},
  pipe title/.style={font=\sffamily\bfseries\fontsize{8}{9.4}\selectfont,
    align=center,text=pipelineInk,text width=2.05cm,inner sep=0pt},
  pipe detail/.style={font=\sffamily\fontsize{7.2}{8.5}\selectfont,
    align=center,text=pipelineInk,text width=2.05cm,inner sep=0pt,anchor=north},
  icon/.style={draw=pipelineBlue,line width=1.15pt},
  icon line/.style={draw=pipelineBlue,line width=.8pt}
]
% The human comparison is the dominant panel, not a minor pipeline step.
\fill[rounded corners=2.5mm,pipelineBlue!4]
  (6.87,-1.83) rectangle (13.53,1.84);
\path[rounded corners=2.5mm,draw=pipelineGreen,
  fill=pipelineMint,line width=.8pt]
  (-.10,-1.83) rectangle (6.68,1.84);
\node[font=\sffamily\bfseries\fontsize{8.6}{10}\selectfont,
  text=pipelineGreenText,inner sep=0pt] at (3.29,1.55)
  {1. Human check: theorem alignment};
\node[font=\sffamily\fontsize{7.5}{9}\selectfont,
  text=pipelineBlue,inner sep=0pt] at (10.20,1.55)
  {LLM proof construction and machine checking};

\node[pipe title] at (.95,1.01) {Paper's\\result};
\node[pipe title,font=\sffamily\bfseries\fontsize{8.8}{10}\selectfont]
  at (5.55,1.01) {QCert\\main theorem};
\node[pipe title] at (7.85,1.01) {2. Formalize\\the proof};
\node[pipe title] at (10.15,1.01) {3. Proof\\checking};
\node[pipe title] at (12.45,1.01) {Faithful\\formalization};

% Source paper, with a circuit and equation.
\begin{scope}[shift={(.95,0)},
  icon/.append style={draw=pipelineGreen},
  icon line/.append style={draw=pipelineGreen}]
  \path[icon,fill=white] (-.35,-.43) -- (.35,-.43) -- (.35,.22)
    -- (.13,.44) -- (-.35,.44) -- cycle;
  \draw[icon line] (.13,.44) -- (.13,.22) -- (.35,.22);
  \draw[icon line] (-.23,.13) -- (.21,.13);
  \draw[icon line] (-.23,-.04) -- (.21,-.04);
  \path[icon line,fill=pipelineMintAccent] (-.09,-.13) rectangle (.09,.04);
  \draw[icon line] (-.23,-.25) -- (.21,-.25);
  \fill[pipelineBlue] (0,-.25) circle (.025);
\end{scope}

% The main theorem is the interface humans compare with the paper.
\begin{scope}[shift={(5.55,0)},
  icon/.append style={draw=pipelineGreen},
  icon line/.append style={draw=pipelineGreen}]
  \path[icon,fill=white,rounded corners=.7mm] (-.39,-.43) rectangle (.39,.38);
  \path[icon line,fill=pipelineMintAccent,rounded corners=.5mm]
    (-.17,.28) rectangle (.17,.45);
  \foreach \b in {.13,-.06,-.25} {
    \draw[icon line] (-.26,\b) -- (-.22,\b-.04) -- (-.14,\b+.05);
    \draw[icon line] (-.04,\b) -- (.25,\b);
  }
\end{scope}

% Proof construction by coding agents.
\begin{scope}[shift={(7.85,0)}]
  \path[icon,fill=white,rounded corners=.7mm] (-.48,-.36) rectangle (.48,.36);
  \draw[icon line] (-.48,.18) -- (.48,.18);
  \foreach \a in {-.35,-.25,-.15} \fill[pipelineBlue] (\a,.27) circle (.02);
  \draw[icon] (-.20,.05) -- (-.31,-.06) -- (-.20,-.17);
  \draw[icon] (.20,.05) -- (.31,-.06) -- (.20,-.17);
  \draw[draw=pipelineYellow,line width=1.25pt] (.06,.06) -- (-.06,-.19);
\end{scope}

% Proof-assistant check.
\begin{scope}[shift={(10.15,0)}]
  \path[icon,fill=white] (0,.46) -- (.39,.29) -- (.34,-.14)
    .. controls (.29,-.31) and (.13,-.41) .. (0,-.48)
    .. controls (-.13,-.41) and (-.29,-.31) .. (-.34,-.14)
    -- (-.39,.29) -- cycle;
  \draw[draw=pipelineYellow,line width=2pt]
    (-.18,-.01) -- (-.04,-.15) -- (.22,.14);
\end{scope}

% Verified theorem.
\begin{scope}[shift={(12.45,0)}]
  \path[icon,fill=white,rounded corners=.6mm] (-.37,-.40) rectangle (.37,.42);
  \node[font=\bfseries\fontsize{14}{14}\selectfont,text=pipelineBlue,
    inner sep=0pt] at (0,.17) {$\vdash$};
  \draw[icon line] (-.22,-.04) -- (.22,-.04);
  \draw[icon line] (-.22,-.18) -- (.10,-.18);
  \path[draw=pipelineBlue,line width=.9pt,fill=pipelineYellowAccent]
    (.26,-.33) circle (.17);
  \draw[icon line] (.18,-.33) -- (.24,-.39) -- (.35,-.27);
\end{scope}

% Humans compare meaning; the proof assistant checks the proof term.
\draw[alignment arrow] (1.63,0) -- (4.87,0);
\node[font=\sffamily\bfseries\fontsize{8}{9.4}\selectfont,
  text=pipelineGreenText,align=center,inner sep=0pt]
  at (3.25,.25) {Same mathematical};
\node[font=\sffamily\bfseries\fontsize{8}{9.4}\selectfont,
  text=pipelineGreenText,align=center,inner sep=0pt]
  at (3.25,-.25) {statement?};
\foreach \a/\b in {6.17/7.23,8.47/9.53,10.77/11.83}
  \draw[pipe arrow] (\a,0) -- (\b,0);

\node[pipe detail] at (.95,-.62) {Claim stated\\in the paper};
\node[pipe detail] at (5.55,-.62) {Theorem stated\\in Lean / Agda};
\node[pipe detail] at (7.85,-.62) {LLM follows the\\paper's argument};
\node[pipe detail,text width=2.2cm] at (10.15,-.62) {Lean / Agda checks\\the formal proof};
\node[pipe detail] at (12.45,-.62) {Aligned claim +\\checked proof};

\node[font=\sffamily\fontsize{7.2}{8.5}\selectfont,
  text=pipelineInk,align=center,inner sep=0pt,text width=6.2cm]
  at (3.29,-1.43) {\textbf{Only manual audit: theorem and definitions}};
\node[font=\sffamily\fontsize{7.5}{9}\selectfont,
  text=pipelineBlue,align=center,inner sep=0pt] at (10.20,-1.43)
  {Generated proofs need no manual audit};

% A small research branch reuses the checked library and returns every
% proposed extension to the human theorem-alignment check.
\node[draw=pipelineYellow,fill=pipelineLemon,
  rounded corners=1.5mm,line width=.7pt,
  minimum width=4.65cm,minimum height=.72cm,inner sep=3pt,
  font=\sffamily\fontsize{7.5}{9}\selectfont,
  align=center,text=pipelineInk] (extensions) at (8.50,-2.48)
  {\textbf{Explore completeness extensions}\\
   Humans + LLMs: new rule sets or circuit families};
\draw[extension arrow] (12.45,-1.83) -- (12.45,-2.48) -- (extensions.east);
\node[font=\sffamily\fontsize{7}{8}\selectfont,
  text=pipelineInk,fill=white,inner sep=1pt]
  at (11.90,-2.28) {reuse};
\draw[extension arrow] (extensions.west) -- (5.55,-2.48) -- (5.55,-1.83);
\end{tikzpicture}
\endgroup

\caption{Using QCert's shared, verified libraries, humans need only check that its main theorem faithfully expresses the paper's result, including the relations, target group, gate matrices, and interpretation. LLM coding agents formalize the source papers' arguments and supply missing proofs or details where needed. Lean or Agda checks the generated proof, so the proof itself needs no human review. The dashed branch shows how humans and LLMs can join forces to explore new completeness results, building on QCert's verified libraries. Each newly obtained result undergoes the same human review and machine checking.}
\label{fig:qcert-pipeline}
\Description{A large mint-green panel highlights the per-paper human check: compare the paper's result with the QCert main theorem and the definitions it refers to. A green double-headed arrow asks whether they express the same mathematical statement. The pipeline continues with an LLM formalizing the paper's arguments and supplying missing proofs or details where needed. Lean or Agda checks the formal proof, yielding a faithful formalization. A lower branch reuses the verified libraries for human-guided, LLM-assisted extensions. Proposed results undergo the same human review of definitions and theorem statements and machine checking of formal proofs.}
\end{figure}


\paragraph{Orientation} In \Cref{sec:background}, we review the two main strategies for proving completeness of a rule set and how minimality is shown. In \Cref{sec:form}, we describe the QCert framework, which adds matrix semantics to
% our earlier library of circuits and presentations~\cite{bian2027popl},
the library of circuits and presentations of~\citet{bian2027popl},
together with our formalization workflow, how QCert keeps completeness from being lost in translation, and a comparison of the Agda and Lean formalizations. In \Cref{sec:qudit}, we present the formalization of qudit Clifford completeness that QCert grew out of, and the mathematical choices it fixed. In \Cref{sec:corpus}, we use QCert to machine-check the completeness and minimality results of sixteen papers. In \Cref{sec:findings}, we describe the issues these formalizations found in the source papers. In \Cref{sec:conclusion}, we discuss the current limitations of QCert and our plans to extend it.


%In \Cref{sec:toffoli-hadamard}, we use QCert to prove a new result, the completeness of Toffoli--Hadamard circuits on any number of qubits. 
\section{Background}\label{sec:background}
A circuit is a word over an alphabet~$\mathcal{G}$ of gates. A rewrite rule is then a relation between words over this alphabet, and a rule set $\mathcal{R}$ generates a congruence on circuits. To a circuit $C$ we associate an \emph{interpretation} $\llbracket C \rrbracket$, which is its underlying matrix. The completeness results we formalize have matrix entries in the finitely generated rings $\Z\left[1/\sqrt2,i\right]$, $\Z\left[1/2,i\right]$, $\Z\left[1/\sqrt2\right]$, and $\Z\left[1/2\right]$, which make them exactly representable in Agda. We leave the formalization of universal circuits taking values in $\mathbb{C}$ for future work.

There are two main strategies for proving completeness of a rule set $\mathcal{R}$ over a gate set $\mathcal{G}$, both of which we support. Consider any two circuits $C_1$ and $C_2$ over $\mathcal{G}$ that have the same linear map, $\llbracket C_1\rrbracket = \llbracket C_2 \rrbracket$. We need to show that $C_1$ can be rewritten into $C_2$ using $\mathcal{R}$. \Cref{fig:normalform} outlines the first strategy. It specifies a unique normal form for each linear map~\cite{lafont2003boolean,sakashita2013cnot,feng2026complete,feng2026ancillas,selinger2015clifford,bian2026qupit,makary2021real,amy2018cnotdihedral,blake2026phaseaffine,li2025qutrit}, often by giving an exact synthesis algorithm~\cite{greylyn2014,bian2021un,li2021on,fang2026hadamard,clement2019thesis}. Completeness follows by proving that every circuit can be transformed into its normal form using the rewrite rules. By uniqueness, $C_1$ and $C_2$ have the same normal form $N$. Reversing the path from $C_2$ to $N$ and concatenating it with the path from $C_1$ to $N$ then gives the desired path of rewrites.
% \[
% C_1\rewrites{\mathcal{R}}N\rewrites{\mathcal{R}}C_2.
% \]

% \vspace{-.2 cm}

\begin{figure}[!thb]
\centering
\begin{subfigure}[b]{0.42\linewidth}
  \centering
  \scalebox{.8}{\tikzfig{figures/normalform}}
  \caption{The rewrites from $C_1$ and $C_2$ to their common normal form $N$ using $\mathcal{R}$ give rewrites from $C_1$ to $C_2$ using $\mathcal{R}$.}
  \label{fig:normalform}
\end{subfigure}\hfill
\begin{subfigure}[b]{0.52\linewidth}
  \centering
  \scalebox{.8}{\tikzfig{figures/translation}}
  \caption{Completeness of $\mathcal{R}'$ gives rewrites between $C_1'=E(C_1)$ and $C_2'=E(C_2)$. Translating this proof back gives rewrites between $C_1$ and $C_2$ using $\mathcal{R}$.}
  \label{fig:translation}
\end{subfigure}
% \captionsetup{singlelinecheck=false}
\vspace{-.2 cm}
\caption{Two strategies for proving completeness of $\mathcal{R}$ over a gate set $\mathcal{G}$. $\;\rewrites{\mathcal{R}}\;$ means applying a finite sequence of rewrites using $\mathcal{R}$ in either direction, where $*$ denotes zero or more steps.}
\label{fig:strategies}
\end{figure}

The exact synthesis results above present groups of matrices rather than circuits (the first block of \Cref{tab:corpus}). Their generators are \emph{multi-level} matrices. For an $m\times m$ matrix $M$ and distinct indices $a_1,\dots,a_m$ in $\{0,\dots,n-1\}$, the \emph{$m$-level} matrix $M_{[a_1,\dots,a_m]}$ is the $n\times n$ matrix that acts as $M$ on the standard basis vectors $e_{a_1},\dots,e_{a_m}$ and as the identity on the others. For example, $(-1)_{[a]}$ negates $e_a$, and $\mathrm{X}_{[a,b]}$ exchanges $e_a$ and $e_b$. Circuits on $k$ qubits have matrices of dimension $n=2^k$, but these groups are defined in every dimension~$n$. We introduce multi-level matrices for two reasons. The gap we find in \Cref{sec:findings} is in a presentation by one- and two-level matrices~\cite{fang2026hadamard}, and our new result in \Cref{sec:toffoli-hadamard} builds on a presentation by one-, two-, and four-level matrices~\cite{li2021on}.

For these groups, the normal forms come from exact synthesis, and the proof that every circuit can be transformed into its normal form has the same structure for each algorithm. Each matrix, or \emph{state}, has a \emph{level}, built from quantities such as the least denominator exponent of a column, and the algorithm reduces each state $M$ to the identity in steps that lower the level. These steps form the normal word $N_M$. The \emph{Main Lemma} says that for every state $M$ and generator $g$, a diagram of relations closes the square formed by $g$ and the algorithm's next steps from $M$ and $gM$, passing only through states of lower level (\Cref{fig:mainlemma}). By well-founded induction on the level, the relations then rewrite $N_{gM}\,g$ into $N_M$ (\Cref{fig:level}), and by induction on length, every word that reduces $M$ to the identity rewrites into $N_M$ (\Cref{fig:length}). Two words with the same matrix therefore rewrite into each other. These inductions have the same form for every presentation, so the main work in each formalization is its Main Lemma.

\begin{figure}[!thb]
\centering
\begin{subfigure}[b]{0.26\linewidth}
  \centering
  \tikzfig{../planqc/figures/mainlemma}
  \caption{The Main Lemma.}
  \label{fig:mainlemma}
\end{subfigure}\hfill
\begin{subfigure}[b]{0.38\linewidth}
  \centering
  \tikzfig{../planqc/figures/levelinduction}
  \caption{Induction on the level.}
  \label{fig:level}
\end{subfigure}\hfill
\begin{subfigure}[b]{0.32\linewidth}
  \centering
  \tikzfig{../planqc/figures/lengthinduction}
  \caption{Induction on length.}
  \label{fig:length}
\end{subfigure}
\vspace{-.2 cm}
\caption{Completeness from the Main Lemma. Plain arrows are words, double arrows are steps or runs of the algorithm, and dashed arrows are words through states of lower level.}
\label{fig:proof}
\end{figure}

\Cref{fig:translation} illustrates the second strategy, which uses an already complete rule set $\mathcal{R}'$ over a different gate set $\mathcal{G}'$~\cite{bian2023cliffordt2,bian2023cliffordcs3,clement2026realcliffordch,amy2024toffoli,blake2026simpler,clement2023complete,clement2024extensions,delorme2026control,schober2026control,blake2026qudit}. Here, $E$ and $D$ translate circuits from $\mathcal{G}$ to $\mathcal{G}'$ and back, replacing each gate with a circuit over the other gate set. Both $E$ and $D$ preserve circuit composition and the underlying linear map. To transfer completeness from $\mathcal{R}'$ to $\mathcal{R}$, we first show that every rule of $\mathcal{R}'$, translated by $D$, can be derived using $\mathcal{R}$. We can then replace each step of the proof between $C_1'$ and $C_2'$ by its translated derivation, giving $D(C_1')\rewrites{\mathcal{R}}D(C_2')$. 
It then remains to prove that the rewrite $C_j\rewrites{\mathcal{R}}D(C_j')$ is derivable, for which it suffices to prove the round-trip rewrites $g\rewrites{\mathcal{R}}D(E(g))$ for each $g\in\mathcal{G}$.
% We also need to check that each gate can be rewritten to its round-trip translation: . Applying these rewrites gate by gate gives $C_j\rewrites{\mathcal{R}}D(C_j')$ for $j\in\{1,2\}$. Combining these paths yields
% \[
% C_1\rewrites{\mathcal{R}}D(C_1')\rewrites{\mathcal{R}}D(C_2')\rewrites{\mathcal{R}}C_2.
% \]
Sometimes the two strategies are combined, as in the scaled dyadic extension of~\cite{li2021on}.
 
Starting from a complete rule set, we can remove a redundant rule by deriving it from the other rules, thereby preserving completeness. To show that the resulting rule set is minimal, we need to give, for each remaining rule $R$, an interpretation of circuits that every other rule preserves but $R$ does not. Any derivation using the other rules preserves this interpretation, so $R$ cannot be derived from them~\cite{clement2024minimal,blake2026simpler}.
 

\section{The QCert Framework}
\label{sec:form}

Our Agda framework QCert states every completeness theorem as a group presentation.
Unlike the letters of an ordinary word, gates act on wires. The gates in $\mathcal{G}$ are indexed by the number of wires they act on. Wires are numbered from $0$ at the bottom, and a gate always acts on the bottom wires of a circuit. To place it anywhere else, we apply a shift, which takes a circuit on $n$ wires and places it on wires $1,\dots,n$ of a circuit with $n+1$ wires. Each relation is stated on as many wires as the paper draws it on. To apply it inside a circuit with more wires, we need \emph{structural rules}, which hold for every circuit, regardless of its gate set. These rules say that an equation can be shifted up by one wire, gates on disjoint wires commute, and a scalar can be placed on any wire. These define the equations of a PROP~\cite{MacLane1965categorical}, and are often defined implicitly in completeness papers, but a formalization in Agda must define them explicitly.

We associate to circuits a matrix semantics $\llbracket - \rrbracket$ with entries in a finitely generated ring like $\Z\left[1/\sqrt2,i\right]$. Hence, $\llbracket - \rrbracket$ defines a group homomorphism to a unitary matrix group $G$. We build $G$ from the group structures of Agda's standard library, which is checked only once per ring. As matrices, a gate on the bottom wires and a shift are Kronecker products with an identity, so the extension of $\llbracket - \rrbracket$ from gates to circuits is defined once for all gate sets. Since the structural rules are universal, their soundness is also proved only once. Then for each paper, it suffices to check the soundness of its own rules.
Soundness and completeness say the relations present the image of $\llbracket - \rrbracket$, which is the subgroup of $G$ generated by the gates. In Agda, mathematical structures such as groups are defined as records.
% Our earlier library~\cite{bian2027popl} expresses this as a record.
The record $\mathcal{R}$~\aid{IsSubPresentationOf}~$G$ contains the semantics $\llbracket - \rrbracket$, with proofs of soundness and completeness. The record $\mathcal{R}$~\aid{IsPresentationOf}~$G$ adds surjectivity, so the relations present all of $G$.


\Cref{fig:qcert-soundness-completeness} summarizes the two obligations: checking each rule's matrix equality for soundness, and deriving a known complete rule set $\mathcal T$ from $\mathcal R$ for completeness.

\begin{figure}[!htbp]
\centering
% Shared semantics and completeness by rule derivation. Native, editable TikZ.
\begingroup
\definecolor{scBlue}{HTML}{205886}
\definecolor{scGreen}{HTML}{32866B}
\definecolor{scGreenText}{HTML}{23674F}
\definecolor{scMint}{HTML}{EFF9F3}
\definecolor{scMintAccent}{HTML}{D1ECDE}
\definecolor{scInk}{HTML}{233447}
\begin{tikzpicture}[
  x=1cm,y=1cm,line cap=round,line join=round,
  sc arrow/.style={-{Stealth[length=1.6mm,width=1.4mm]},
    draw=scBlue,line width=.85pt},
  sc sound/.style={-{Stealth[length=1.7mm,width=1.5mm]},
    draw=scGreen,line width=1.1pt},
  sc card/.style={draw=scBlue!30,fill=white,rounded corners=1.5mm,
    line width=.6pt},
  sc title/.style={font=\sffamily\bfseries\fontsize{9}{10.5}\selectfont,
    text=scInk,align=center,inner sep=0pt,text width=5.7cm},
  sc text/.style={font=\sffamily\fontsize{7.7}{9.1}\selectfont,
    text=scInk,align=center,inner sep=0pt,text width=5.65cm},
  sc math/.style={font=\fontsize{10}{11.5}\selectfont,
    text=scInk,inner sep=0pt}
]
% The two proof obligations share a single matrix interpretation.
\fill[scMint,rounded corners=2.5mm] (0,.75) rectangle (6.22,7.85);
\fill[scBlue!4,rounded corners=2.5mm] (7.52,.75) rectangle (13.74,7.85);
\node[sc title,text=scGreenText] at (3.11,7.52)
  {(a) Soundness of $\mathcal R$};
\node[sc title,text=scBlue] at (10.63,7.52)
  {(b) Completeness of $\mathcal R$};

% The generator-level definition is the most prominent card.
\path[draw=scGreen,fill=scMintAccent,rounded corners=1.8mm,line width=1.1pt]
  (.20,5.25) rectangle (6.02,7.10);
\node[sc title,text=scGreenText] at (3.11,6.81)
  {1. Interpret generators};
\node[sc text] at (3.11,6.43) {Gates or two-level operators};
\node[sc math,font=\fontsize{12}{13}\selectfont] at (3.11,6.00)
  {$g\longmapsto M_g\in U_n(\D[\omega])$};
\node[sc text,font=\sffamily\bfseries\fontsize{7.7}{9.1}\selectfont,
  text=scGreenText] at (3.11,5.55)
  {Uniform matrix semantics};
\draw[sc arrow] (3.11,5.25) -- (3.11,4.95);

\path[sc card] (.20,3.25) rectangle (6.02,4.95);
\node[sc title] at (3.11,4.62) {2. Lift to circuits / words};
\node[sc math] at (3.11,4.18)
  {$g_1\cdots g_k\quad\longmapsto\quad M_{g_1}\cdots M_{g_k}$};
\node[sc text] at (3.11,3.65)
  {Generic lifting handles\\composition and wire placement.};
\draw[sc arrow] (3.11,3.25) -- (3.11,2.95);

\path[sc card] (.20,1.00) rectangle (6.02,2.95);
\node[sc title] at (3.11,2.62) {3. Compare both sides of every rule};
\node[sc text] at (3.11,2.25) {For each $(u=v)\in\mathcal R$, check};
\node[sc math,font=\fontsize{12}{13}\selectfont] at (3.11,1.80)
  {$\llbracket u\rrbracket=\llbracket v\rrbracket$};
\node[sc text,text=scGreenText,
  font=\sffamily\bfseries\fontsize{7.7}{9.1}\selectfont] at (3.11,1.28)
  {Every $\mathcal R$-rewrite preserves its matrix.};

% A complete reference set may be redundant, but is in the same syntax.
\path[sc card] (7.72,5.25) rectangle (13.54,7.10);
\node[sc title] at (10.63,6.81) {1. Fix a complete reference set};
\node[sc math,font=\fontsize{12}{13}\selectfont] at (10.63,6.28)
  {$\mathcal T=\{A=B,\ldots\}$};
\node[sc text] at (10.63,5.72)
  {Known complete, possibly redundant;\\same language and matrix semantics as $\mathcal R$.};
\draw[sc arrow] (10.63,5.25) -- (10.63,4.95);

\path[sc card] (7.72,2.60) rectangle (13.54,4.95);
\node[sc title] at (10.63,4.62) {2. Derive every reference equation};
\node[sc text] at (10.63,4.20)
  {For every $(A=B)\in\mathcal T$, construct};
\node[sc math] at (8.40,3.62) {$A=C_0$};
\node[sc math] at (9.89,3.62) {$C_1$};
\node[sc math] at (10.97,3.62) {$\cdots$};
\node[sc math] at (12.53,3.62) {$C_k=B$};
\foreach \a/\b in {9.00/9.48,10.19/10.59,11.34/11.91} {
  \draw[{Stealth[length=1.1mm]}-{Stealth[length=1.1mm]},
    draw=scGreen,line width=.8pt] (\a,3.62) -- (\b,3.62);
  \node[font=\fontsize{7.5}{8.5}\selectfont,text=scGreenText,inner sep=0pt]
    at ({(\a+\b)/2},3.89) {$\mathcal R$};
}
% Soundness of R has already been established by the left column.
\node[sc text,text=scGreenText,fill=scMint,rounded corners=1mm,
  inner sep=3pt,text width=5.10cm,
  font=\sffamily\bfseries\fontsize{7.7}{9.1}\selectfont]
  (sound-steps) at (10.63,3.08)
  {Each step uses a sound rule from $\mathcal R$.};
\draw[sc arrow] (10.63,2.60) -- (10.63,2.30);

\path[sc card] (7.72,1.00) rectangle (13.54,2.30);
\node[sc title] at (10.63,1.99) {3. Conclude $\mathcal R$ is complete};
\node[sc text] at (10.63,1.43)
  {\mbox{Every valid equation is derivable from $\mathcal R$.}};

% Soundness certifies the chain; it does not construct or find it.
\draw[sc sound] (6.02,1.80) -- (6.87,1.80) -- (6.87,3.08) -- (sound-steps.west);
\end{tikzpicture}
\endgroup

\caption{QCert's shared matrix interpretation supports both proof obligations. The left column checks every rule of $\mathcal R$; the right derives every rule of a known complete set $\mathcal T$ in the same circuit language and matrix semantics. Each rewrite step uses a rule of $\mathcal R$, whose soundness is checked on the left. The displayed $U_n(\D[\omega])$ covers qubit fragments~\cite{amy2020numbertheoretic}, with $\D=\Z[1/2]$ and $\omega=e^{i\pi/4}$. Qudit circuit families use other rings.}
\label{fig:qcert-soundness-completeness}
\Description{Two columns show soundness and completeness. On the left, the highlighted uniform matrix semantics assigns a matrix to each gate or two-level operator. QCert lifts this assignment to circuits and words, then checks equality of the matrices on both sides of each rule. On the right, every equation in a known complete reference rule set is derived by rewrites from the proposed rule set. The rewrite arrows are green. An arrow from the completed soundness checks points to the statement that each step uses a sound rule from the proposed rule set. Deriving every reference rule establishes completeness.}
\end{figure}

Each formalization ends with its main theorem stated as an \aid{IsSubPresentationOf} or \aid{IsPresentationOf} record, in one format for all of them. The relations are written \aid{n CRel,{\_}==={\_}} for circuits on $n$ wires and \aid{n MRel,{\_}==={\_}} for words over multi-level matrices of dimension $n$. For example, \Cref{agda:stab} is the theorem for the real stabilizer circuits of~\cite{makary2021real}, with the relations R1--R16 of its Figure~1 lifted to $n$ wires. \Cref{agda:hpi} is the theorem for the two-level matrices of Hadamard-$\Pi$~\cite{fang2026hadamard}, with the relations of its Figure~6. Completeness and soundness are read off the record in one form too, \Cref{agda:complete,agda:sound}, where \aid{⊢ w ≈ v} says that $w$ and $v$ are equal by the relations. The qubit developments compare matrices over $\D[\omega]$ (\aid{{\_}≈ᵁ{\_}}), in which each reads its group; the qudit ones compare in their own groups.
\begin{agdacode}[escapechar=@]
presentation : ∀ n → (n CRel,_===_) IsSubPresentationOf U (2 ^ n)@\agdatag{agda:stab}@
presentation : ∀ n → (n MRel,_===_) IsPresentationOf (O n)@\agdatag{agda:hpi}@
complete : ∀ {n} {w v : Circuit n} → ⟦ w ⟧ ≈ᵁ ⟦ v ⟧ → ⊢ w ≈ v@\agdatag{agda:complete}@
sound : ∀ {n} {w v : Circuit n} → ⊢ w ≈ v → ⟦ w ⟧ ≈ᵁ ⟦ v ⟧@\agdatag{agda:sound}@
\end{agdacode}
The minimality results of~\cite{blake2026simpler} take the same form, \Cref{agda:minimal}, for the rule set of a fragment on $n$ wires, \aid{n ⊢{\_}==={\_}}. It is complete if it derives every equation that the matrices validate, and minimal if, without any one of its rules on $n$ wires, it is no longer complete there. Minimality only makes sense for a rule set, which says what dropping a rule is, so Agda accepts \aid{IsMinimal} only for \aid{n ⊢{\_}==={\_}}.
\begin{agdacode}[escapechar=@]
minimal-complete : ∀ {n} → (n ⊢_===_) IsMinimal × (n ⊢_===_) IsComplete@\agdatag{agda:minimal}@
\end{agdacode}

Agda checks every proof, so a reader only needs to check that each theorem statement matches its paper's result. Since every theorem is one of these two records, it suffices for a reader to check only four components: the relations transcribed from the paper, the target group, the matrix of each gate, and the line assigning $\llbracket - \rrbracket$ to the record. Each Agda development collects these components and its main theorem in one module, \aid{MainTheorem}, and the file \texttt{MainTheorems.agda} imports one such module per source and restates each theorem; the corresponding Lean statements are collected in \texttt{QCert/Table1.lean}.

\paragraph{Extensions of the earlier framework} To support these formalizations, QCert extends the Agda framework of~\citet{bian2027popl}. Besides the matrix semantics above and the statement of minimality (\Cref{sec:background}), it adds three features.
\begin{itemize}
\item \emph{Three-wire gates.} The earlier framework allows gates on one and two wires. QCert also allows gates on three wires, such as the Toffoli gate, together with the structural rule they need.
\item \emph{Scalars.} QCert also allows scalars, which act on no wire, such as the global phase $\omega$ of Clifford circuits, with a structural rule that places a scalar on any wire. A rule set can then relate circuits exactly rather than up to a global phase.
\item \emph{PROPs.} In the earlier framework, a swap is an ordinary gate whose laws must be listed among the rules. QCert can also treat circuits as a PROP. One two-wire gate is designated as the swap, and the framework supplies its laws, that it is an involution and that every gate can be moved across it, once for all gate sets. A paper's rule set then lists only its own rules, each on the wires it is drawn on, while the swap and its laws stay hidden. We use this for the circuits of~\cite{amy2018cnotdihedral,blake2026phaseaffine,blake2026simpler}.
\end{itemize}
Where the examples of~\citet{bian2027popl} are groups such as symmetric, cyclic and Pauli groups, QCert formalizes sixteen published completeness and minimality results (\Cref{tab:corpus}).

\subsection{Formalization Workflow}\label{sec:method}

Our workflow has three parts. First, before formalizing a paper, we check its rewrite rules, its intermediate results, and each line of its derivations with exact matrices in short scripts. This caught most typos before any Agda code was written. Second, we state each main theorem in QCert before a coding agent starts, so the agent only supplies the proof. Finally, for long derivations, the agent writes a script that turns each line of the paper into one step of the Agda proof, so a failing step points to a specific line of the paper. 
%With this workflow, coding agents wrote the fifteen formalizations other than qudit Clifford within four weeks, most in less than a day.

The main computational cost came from Agda's conversion checker. If two terms are syntactically identical (i.e., written in the same way), it accepts them immediately. Otherwise, it evaluates both, which can take exponential time for a long matrix product. For such products, we keep circuit interpretations abstract and decide matrix equality by computing a single Boolean using row operations, which is orders of magnitude faster than comparing matrices entry by entry. We record these techniques in the instructions followed by the agents, so later proofs avoid the same slowdowns.
\subsection{Avoiding Completeness Lost in Translation}
\label{sec:lost}

A machine-checked proof establishes only what its statement says. As the Navier--Stokes example in the introduction shows, a formalization can end up verifying a weaker result than the one intended~\cite{bastounis2026navier}. For a completeness result, the meaning can be lost when the rules, the semantics, or the theorem is translated into a proof assistant. QCert addresses this in three ways.

\paragraph{Foundations} QCert builds on our formalization of qudit Clifford completeness and keeps the two mathematical choices that \Cref{sec:qudit} explains. The semantics is constructed independently of the circuit syntax. The target group is defined by its own mathematics, as a group of matrices or a semidirect product, and not as a quotient of circuits, so completeness relates two independently given objects and cannot hold by construction. Each formalization also states what it does and does not verify. LLMs help construct the proofs and the proof assistant checks them, but we remain responsible for the definitions, the proof strategy, and the connection to the intended mathematics.

\paragraph{Explicit proof obligations} A proof by the second strategy of \Cref{sec:background} starts from the sound rules $\mathcal{R}$ of the source paper and uses them to derive the translation of every rule of a rule set known to be complete (\Cref{fig:translation}). This strategy suits QCert, because it turns completeness into explicit proof obligations, one for each rule of the complete set, which Agda either accepts or rejects. When an obligation cannot be discharged, say a rule $r$ does not follow from $\mathcal{R}$, there are two possibilities. Either a derivation or a case is missing, and we supply it, or $r$ must be added to $\mathcal{R}$. In the second case, the completeness claim concerns the revised rule set, and the formal statement says so. The relation $r_{\mathrm K}$ of \Cref{sec:findings} is an example. It does not follow from the rules of~\cite{amy2024toffoli}, so we added it and proved that it is independent of the others.

\paragraph{Shared structure} The matrix groups of different circuit families are mathematically related, often as subgroups of one another. The first five rows of \Cref{tab:corpus} present five of these groups, and several circuit results are proved by translating into them. For example, two-qubit Clifford+T circuits are translated into $U_4\left(\Z\left[1/\sqrt2,i\right]\right)$, three-qubit Clifford+CS circuits into $U_8\left(\Z\left[1/2,i\right]\right)$, and real Clifford+CH circuits into $O_{2^n}\left(\Z\left[1/\sqrt2\right]\right)$. A common framework reflects this structure rather than merely collecting separate proofs. We verify general arguments and algebraic constructions once, such as Reidemeister--Schreier coset enumeration and presentations of semidirect products, and reuse them by checking only the assumptions and semantic connections specific to each family. This reduces the need to translate each proof independently and makes the dependencies of each completeness claim explicit.

\subsection{Comparison of the Agda and Lean Formalizations}
\label{sec:compare}

The Lean~\cite{moura2021lean4} version of QCert formalizes the same sixteen results on top of Mathlib~\cite{mathlib2020}, which makes the corpus available in the Lean ecosystem. Its framework is not an independent implementation. We instructed the coding agent that wrote it not to consult our Agda implementation, but the Lean framework is nonetheless a translation of the Agda one. Its words, circuit generators, structural rules, and presentation and sub-presentation predicates are those of the Agda framework, down to the names of the definitions and the wording of their documentation. The two formalizations therefore share the design of their framework, and would share any mistake in it, so their agreement is weaker evidence than two independent formalizations would give. They use different representations and sometimes different proofs, so the theorem statements of each must be checked against the source papers separately. Below we compare their definitions, and then their proof development, verification time, and execution time (\Cref{fig:compare} and \Cref{tab:compare}). These observations concern our two formalizations, not Agda and Lean in general.

\paragraph{Definitions and semantics} The translation keeps the definitions but adapts their representation to Mathlib. Agda uses its own word type, with explicit brackets, and the setoid-based algebra of the standard library. A presentation is a record containing a semantics and a group isomorphism, or a monomorphism for a sub-presentation. Lean uses Mathlib's free monoid and the monoid congruence generated by the relations. Its presentations are predicates in \texttt{Prop}, from which it derives isomorphisms out of the quotient. In both, the circuit infrastructure, the structural rules, and general completeness arguments are shared between developments.

The semantic targets also differ. Most Agda developments work directly with exact representations of matrix entries. Lean states its matrix groups over Mathlib's real or complex numbers and proves bridges to them from computable exact rings. In particular, the Agda qudit Clifford development presents the constructed model of \Cref{sec:qudit}, whereas Lean proves soundness and completeness against complex matrices for every odd prime. A common source therefore does not imply identical statements or intermediate lemmas.

\begin{figure}[tbp]
\centering
% Generated by pldi/compare/mkfig.py from the measurements of
% pldi/compare/ (code lines: loc.py; check times: Table 1).
\begin{tikzpicture}
\begin{groupplot}[
  group style={group size=2 by 1, horizontal sep=18pt,
               yticklabels at=edge left},
  width=0.40\linewidth, height=9.2cm, scale only axis=false,
  y dir=reverse, ymin=-0.7, ymax=16.7,
  ytick={0,...,16}, yticklabels={{$U_n(\Z[1/\sqrt2,i])$},{$U_n(\Z[1/2,i])$},{$O_n(\Z[1/2])$, $L_n$},{$O_n(\Z[1/\sqrt2])$, Hadamard-$\Pi$},{$O_n(\Z[1/\sqrt2])$, Cl\'ement},{Clifford},{Clifford, odd prime $p$},{real stabilizer},{CNOT-dihedral},{phase-affine, prime $d$},{two-qubit Clifford+T},{three-qubit Clifford+CS},{real Clifford+CH},{Toffoli--Hadamard},{qutrit Clifford},{minimality, six fragments},{shared library}},
  y tick label style={font=\scriptsize},
  x tick label style={font=\scriptsize},
  xlabel style={font=\footnotesize},
  title style={font=\footnotesize\bfseries, yshift=-2pt},
  tick align=outside, tick pos=lower,
  axis line style={cmpAxis}, major tick style={cmpAxis},
  xmajorgrids, grid style={cmpGrid},
  clip=false,
]
% (a) Proof development: code lines, written and generated.
\nextgroupplot[
  title={(a) Proof development},
  xbar, bar width=3.1pt,
  xmin=0, xmax=140, xtick={0,20,...,120},
  xlabel={Code lines (thousands)},
  legend to name=cmpLegend, legend columns=4,
  legend style={font=\scriptsize, draw=none, column sep=6pt},
  legend image code/.code={\draw[#1] (0pt,-2.5pt) rectangle (9pt,2.5pt);},
]
% Each system's total, hatched, with its written part drawn over it.
\addplot[forget plot, bar shift=1.9pt, draw=cmpAgda, pattern=north east lines,
         pattern color=cmpAgda]
  coordinates {(7.21,0) (7.71,1) (13.58,2) (9.73,3) (9.09,4) (10.75,5) (57.83,6) (11.06,7) (5.04,8) (10.73,9) (11.03,10) (56.24,11) (60.17,12) (28.57,13) (53.85,14) (23.09,15) (13.12,16)};
\addplot[forget plot, bar shift=1.9pt, fill=cmpAgda, draw=cmpAgda]
  coordinates {(7.21,0) (7.71,1) (11.87,2) (7.64,3) (7.14,4) (3.89,5) (57.6,6) (5.1,7) (4.96,8) (10.59,9) (0.63,10) (0.72,11) (58.57,12) (6,13) (4.46,14) (8.62,15) (10.83,16)};
\addplot[forget plot, bar shift=-1.9pt, draw=cmpLean, pattern=north west lines,
         pattern color=cmpLean]
  coordinates {(4.43,0) (3.92,1) (9.85,2) (13.51,3) (2.43,4) (5.43,5) (6.37,6) (5.04,7) (3.65,8) (7.02,9) (5.88,10) (138.85,11) (105.89,12) (3.74,13) (10.21,14) (13.37,15) (1.22,16)};
\addplot[forget plot, bar shift=-1.9pt, fill=cmpLean, draw=cmpLean]
  coordinates {(4.43,0) (3.92,1) (6.2,2) (6.65,3) (1.18,4) (4.97,5) (5.95,6) (2.41,7) (2.75,8) (6.24,9) (0.76,10) (1.53,11) (44.15,12) (1.19,13) (5.14,14) (5.52,15) (1.12,16)};
\addlegendimage{fill=cmpAgda, draw=cmpAgda}
\addlegendentry{Agda, written}
\addlegendimage{draw=cmpAgda, pattern=north east lines, pattern color=cmpAgda}
\addlegendentry{Agda, generated or ported}
\addlegendimage{fill=cmpLean, draw=cmpLean}
\addlegendentry{Lean, written}
\addlegendimage{draw=cmpLean, pattern=north west lines, pattern color=cmpLean}
\addlegendentry{Lean, generated}
\draw[cmpAxis, densely dotted] (axis cs:0,15.5) -- (axis cs:140,15.5);
% (b) Verification time: check time per row, log scale.
\nextgroupplot[
  title={(b) Verification time},
  xmode=log, xmin=10, xmax=20000, log basis x=10,
  xlabel={Check time (s), log scale},
]
\addplot[cmpLink, forget plot] coordinates {(42,0) (115.62,0)};
\addplot[cmpLink, forget plot] coordinates {(62,1) (140.99,1)};
\addplot[cmpLink, forget plot] coordinates {(79,2) (157.76,2)};
\addplot[cmpLink, forget plot] coordinates {(44,3) (315.7,3)};
\addplot[cmpLink, forget plot] coordinates {(64,4) (52.94,4)};
\addplot[cmpLink, forget plot] coordinates {(34,5) (103.61,5)};
\addplot[cmpLink, forget plot] coordinates {(546,6) (166.92,6)};
\addplot[cmpLink, forget plot] coordinates {(35,7) (62.87,7)};
\addplot[cmpLink, forget plot] coordinates {(24,8) (77.29,8)};
\addplot[cmpLink, forget plot] coordinates {(61,9) (176.59,9)};
\addplot[cmpLink, forget plot] coordinates {(118,10) (275.4,10)};
\addplot[cmpLink, forget plot] coordinates {(250,11) (1162.66,11)};
\addplot[cmpLink, forget plot] coordinates {(405,12) (8868.59,12)};
\addplot[cmpLink, forget plot] coordinates {(102,13) (328.84,13)};
\addplot[cmpLink, forget plot] coordinates {(282,14) (164.76,14)};
\addplot[cmpLink, forget plot] coordinates {(183,15) (300.2,15)};
\addplot[only marks, mark=*, mark size=1.9pt, color=cmpAgda]
  coordinates {(42,0) (62,1) (79,2) (44,3) (64,4) (34,5) (546,6) (35,7) (24,8) (61,9) (118,10) (250,11) (405,12) (102,13) (282,14) (183,15)};
\addplot[only marks, mark=square*, mark size=1.7pt, color=cmpLean]
  coordinates {(115.62,0) (140.99,1) (157.76,2) (315.7,3) (52.94,4) (103.61,5) (166.92,6) (62.87,7) (77.29,8) (176.59,9) (275.4,10) (1162.66,11) (8868.59,12) (328.84,13) (164.76,14) (300.2,15)};
\node[font=\scriptsize, text=black!75, anchor=west, xshift=3pt]
  at (axis cs:546,6) {Agda};
\node[font=\scriptsize, text=black!75, anchor=east, xshift=-3pt]
  at (axis cs:166.92,6) {Lean};
\end{groupplot}
\node[anchor=north] at ($(group c1r1.south east)!0.5!(group c2r1.south west) + (0,-28pt)$)
  {\pgfplotslegendfromname{cmpLegend}};
\end{tikzpicture}

\par\medskip
\begin{tikzpicture}
\begin{axis}[
  title={(c) Execution time},
  title style={font=\footnotesize\bfseries, yshift=-2pt},
  width=0.62\linewidth, height=5.2cm,
  xmin=1.7, xmax=8.3, xtick={2,...,8},
  ymode=log, log basis y=10, ymin=1, ymax=100000,
  xlabel={Qubits}, ylabel={Time per decision (ms), log scale},
  xlabel style={font=\footnotesize}, ylabel style={font=\footnotesize},
  tick label style={font=\scriptsize},
  tick align=outside, tick pos=lower,
  axis line style={cmpAxis}, major tick style={cmpAxis},
  ymajorgrids, grid style={cmpGrid}, clip=false,
]
\addplot[color=cmpAgda, line width=1pt, mark=*, mark size=1.9pt]
  coordinates {(2,5) (3,39.6) (4,335.2) (5,2765.8) (6,24356.6)};
\addplot[color=cmpLean, line width=1pt, mark=square*, mark size=1.7pt]
  coordinates {(2,5.8) (3,18.8) (4,56.4) (5,206.6) (6,612) (7,2314.8) (8,7895.4)};
\node[font=\scriptsize, text=black!75, anchor=west, xshift=4pt]
  at (axis cs:6,24356.6) {Agda};
\node[font=\scriptsize, text=black!75, anchor=west, xshift=4pt]
  at (axis cs:8,7895.4) {Lean};
\end{axis}
\end{tikzpicture}

\caption{The Agda and Lean formalizations of each row of \Cref{tab:corpus}. (a) Code lines, excluding blank lines and comments. The last row is the shared library of each framework. (b) The check times of \Cref{tab:corpus}, Agda (circles) and Lean (squares), measured under the different conditions described in the text. (c) The time each compiled formalization takes to decide that two Clifford circuits of 100 gates have the same matrix, by the procedure its soundness proof uses, as medians over three runs of five random circuits.}
\label{fig:compare}
\end{figure}

\paragraph{Proof development} Excluding blank lines and comments, the Agda formalization has about 376{,}000 lines and the Lean formalization about 339{,}000 (\Cref{fig:compare}(a)). Lean's is shorter in thirteen of the sixteen rows, by a factor of five or more for qudit Clifford, Toffoli--Hadamard, and qutrit Clifford. It is longer for Hadamard-$\Pi$, three-qubit Clifford+CS, and real Clifford+CH. Both counts include the derivations and certificates that scripts generate. The Agda proofs are explicit terms, built by pattern matching and equational reasoning, with reflection-based solvers for associativity and ring identities. Its library develops the group theory the proofs need, in about 13{,}100 lines against 1{,}200 in Lean. The Lean proofs rely on Mathlib's algebra and on tactics, with about 69{,}000 calls to automation such as \texttt{simp}, \texttt{omega}, \texttt{decide}, and \texttt{ring}.

Both implementations reuse normal-form, exact-synthesis, and translation arguments. Both also check generated derivations with a replay procedure that is proved correct. The proof assistant checks the computation and the theorem that a successful replay gives a derivation, so the scripts that search for derivations need not be trusted. For Clifford+CS, Lean constructs its own coset certificates. The trust checks reflect the different foundations. Every Agda module is \texttt{-{}-safe} and uses no postulates, so the Agda corpus is constructive. The Lean corpus uses no \texttt{native\_decide}, and its main theorems depend on no unfinished proof. An audit of 105 declarations reports only the axioms \texttt{propext}, \texttt{Classical.choice}, and \texttt{Quot.sound}.

\paragraph{Verification time} \Cref{tab:corpus} and \Cref{fig:compare}(b) give the check times. In total Agda took 2{,}332 seconds and Lean 9{,}305, or 1{,}927 and 3{,}826 without real Clifford+CH, whose Lean build alone took 5{,}479 seconds, more than half of the total. Its five-wire certificate modules \texttt{Family267} and \texttt{Family268} alone took about 1{,}741 and 1{,}369 seconds. Lean was faster for three rows, Cl\'ement's presentation, qudit Clifford, and qutrit Clifford, whose Lean proofs take different routes. Both assistants decide equalities of matrices and certificates by evaluation in the checker. Agda reduces a Boolean to \texttt{true}, and Lean uses \texttt{decide +kernel}.

The two columns are not a matched comparison. The Agda column keeps the recorded times of \Cref{tab:corpus}, which were taken on a different machine. The Lean column comes from a clean build of revision \texttt{e19af70c} on 5--6~November~2026, with Lean~4.30.0 and the pinned Mathlib revision, under WSL on an Intel Core i5-13400TEF with sixteen logical CPUs, about 26\,GiB of RAM, and 64\,GiB of swap. We removed QCert's build artifacts, kept the Mathlib and other dependency caches, and built the folders one at a time in dependency order with \texttt{LEAN\_NUM\_THREADS=2}. Each time is the elapsed wall time of a folder, including Lake startup and artifact generation, and a shared dependency is counted only in the folder that defines it. These are therefore build times and not kernel CPU times. Proof constructions, module boundaries, parallelism, and artifact generation also differ, so the differences cannot be attributed to the proof assistants alone. We have yet to separate the time Lean spends running automation from the time it spends checking the resulting proofs.

The build covered all 810 tracked library modules, excluding three generator templates, and took \textbf{9{,}415.37 seconds (2\,h 36\,min 55\,s)} in timed build commands. Of these modules, 808 passed. Two unused legacy helpers in real Clifford+CH failed, \texttt{ItemBFun} at \texttt{omega} proofs and \texttt{ItemASteps} at the heartbeat limit. Neither is imported by the main theorem chain or the audit, both of which passed, and the real Clifford+CH entry of \Cref{tab:corpus} includes the whole attempted build. Beyond the sixteen development folders, the shared \texttt{Circuit}, \texttt{Exact}, \texttt{Matrix}, and \texttt{NumberTheory} folders took 17.91, 15.51, 6.88, and 15.29 seconds, and the five root modules, including the audit, took 54.95 seconds. Lean builds of another project ran on the same machine during part of the run. We therefore re-timed the twenty folders and root modules they overlapped and report the second timing, during which the other project used about a quarter of one CPU on average. The two runs total 9{,}604 and 9{,}415 seconds, and the two timings of each development folder agree within 15\%, except for Toffoli--Hadamard (178 and 135 seconds).
% Measurement record: pldi/compare/leantime.json (written by pldi/compare/leantime.py).

\paragraph{Execution time} The Agda normal forms are executable functions. The Lean developments obtain normal forms either as existence theorems or as \texttt{noncomputable} definitions. For example, the Lean proofs for the Clifford and real stabilizer circuits of~\cite{selinger2015clifford,makary2021real} show that every circuit equals a normal form without defining a function that computes it, and the exact synthesis of the two-level groups is noncomputable. In all, 460 definitions in the sixteen Lean developments are declared \texttt{noncomputable}, so there is no Lean normalizer to time. Both formalizations can, however, execute the procedure that decides whether two circuits have the same matrix, which their soundness proofs run. We compiled this procedure for the Clifford circuits of~\cite{selinger2015clifford} to native code, through Agda's GHC backend and Lean's C backend. We then timed both on an Intel Core i5-14400, on the same random circuits of 100 gates over $H$, $S$, and $CZ$, deciding that each circuit $w$ has the matrix of $w\,S^4$ (\Cref{fig:compare}(c)). On two qubits Agda is slightly faster, at 5.0 against 5.8 milliseconds per decision. From three qubits on Lean is faster, by a factor of 2.1 on three qubits and 40 on six, where Agda takes 24 seconds and Lean 0.61. The gap comes mostly from the procedures rather than the compilers. Agda computes each circuit's matrix as a product of stored $2^n\times 2^n$ matrices, so its work grows about eightfold with each qubit. Lean simulates the circuit on each basis vector, so its work grows about fourfold.

\begin{table}[tbp]
\small
\setlength{\tabcolsep}{4pt}
\renewcommand{\arraystretch}{1.3}
\begin{tabularx}{\linewidth}{@{}l Y Y@{}}
\toprule
& Agda & Lean \\
\midrule
\multicolumn{3}{@{}l}{\emph{Proof development}}\\
Logic & constructive: every module \texttt{-{}-safe}, no postulates & classical: every main theorem uses \texttt{Classical.choice} \\
Library & standard library (104 modules imported), EucDomain & Mathlib (56 modules imported directly) \\
Proofs & terms, by pattern matching and equational reasoning; reflection-based solvers & tactics: about 75{,}000 tactic blocks, 69{,}000 automation calls \\
Reused from the sources & the authors' Agda proofs for Clifford+T and Clifford+CS & no proofs; the authors' coset-action data for Clifford+T \\
Modules & 1{,}476 & 795 \\
Code lines & 375{,}647 & 338{,}866 \\
\midrule
\multicolumn{3}{@{}l}{\emph{Verification time}}\\
Decided equalities & a Boolean reduced to \texttt{true} & \texttt{decide +kernel} \\
Check time, sixteen rows & 2{,}332\,s & 9{,}305\,s \\
Faster for & thirteen rows & three rows \\
\midrule
\multicolumn{3}{@{}l}{\emph{Execution time}}\\
Normal forms & executable functions & existence theorems or \texttt{noncomputable} definitions (460 declared) \\
Benchmarked procedure & matrix of each circuit, by products of stored matrices & simulation of each circuit on every basis vector \\
Compiled through & GHC (Haskell) & C \\
Per decision, Clifford, 100 gates & 5.0\,ms on 2 qubits, 24\,s on 6 & 5.8\,ms on 2 qubits, 0.61\,s on 6 \\
\bottomrule
\end{tabularx}
\caption{The Agda and Lean formalizations compared. Counts are over the files under version control in the sixteen development folders of \Cref{tab:corpus}, without the shared libraries, and code lines exclude blank lines and comments. Check times are those of \Cref{tab:corpus}, and execution times those of \Cref{fig:compare}(c).}
\label{tab:compare}
\end{table}

\section{The Foundation: Qudit Clifford Completeness}
\label{sec:qudit}

QCert grew out of our formalization of the completeness theorem of~\quditcite{}. That theorem gives sixteen rewrite rules that are sound and complete for multi-qudit Clifford circuits over $H$, $S$, and $CZ$ in every odd prime dimension~$p$. Unlike the other fifteen formalizations, this one was built largely by hand, and it fixed the mathematical choices that the rest of QCert follows. It is checked in safe, axiom-free Agda against a semantics constructed independently of the syntax. The proof follows the paper's except at one step, where we instead \emph{factor the group, present the factors, and compose the presentations} (\Cref{sec:q-circuits}).

\subsection{Qudit Clifford Circuits and the Main Theorems}\label{sec:q-circuits}

We use the gate conventions of~\quditcite{}. Fix an odd prime $p$ and let $\omega = e^{2\pi i/p}$. On the basis $\ket{0},\dots,\ket{p-1}$ indexed by $\Zp$, the gates $X\ket{j} = \ket{j+1}$ and $Z\ket{j} = \omega^{j}\ket{j}$ generate, with the phases $\omega^t$, the Pauli group $\Pauli$ on $n$ qudits. The gates
\[
H\colon \ket{j} \mapsto \tfrac{\delta_p}{\sqrt p}\sum_k \omega^{jk}\ket{k}, \qquad
S\colon \ket{j} \mapsto \omega^{j(j-1)/2}\ket{j}, \qquad
CZ\colon \ket{j}\ket{l} \mapsto \omega^{jl}\ket{j}\ket{l}
\]
generate, with the same phases, the Clifford group $\Cliff$. Here $\delta_p = (-i)^{(p-1)/2}$, which makes the determinant of $H$ equal to~$1$. Modulo scalars, we write $\PCliff = \Cliff/\langle\omega\rangle$ and $\PPauli = \Pauli/\langle\omega\rangle \cong \Zp^{2n}$. Then $\Cliff/\Pauli \cong \PCliff/\PPauli \cong \Sp$ and $\PCliff \cong \PPauli \rtimes \Sp$. We do not formalize these last two isomorphisms, because our semantic groups are \emph{defined} as their right-hand sides. \quditcite{} also include $-1$ among the global phases, so that the swap and the multiplier gates belong to $\Cliff$. We work mainly in the projective quotient, where $-1$ makes no difference and the relations are slightly simpler.

Using the framework of \Cref{sec:form}, we encode circuits, rules, and reasoning with the rules. We draw circuits in matrix-multiplication order, the reverse of the usual circuit order, so $w \bullet v$ applies $v$ first but draws it on the right. For example, this is the swap:

\noindent\begin{minipage}[c]{0.56\textwidth}
%% Generated by make agda from lagda/ex-circuit.lagda.tex (agda --latex --only-scope-checking);
%% edit that file, not this one.  The hidden block renders nothing.
\begin{code}[hide]%
\>[0]\AgdaKeyword{module}\AgdaSpace{}%
\AgdaModule{vqupit-short.lagda.ex-circuit}\AgdaSpace{}%
\AgdaKeyword{where}\<%
\\
%
\\[\AgdaEmptyExtraSkip]%
\>[0]\AgdaKeyword{open}\AgdaSpace{}%
\AgdaKeyword{import}\AgdaSpace{}%
\AgdaModule{vqupit-short.lagda.Prelude}\<%
\\
%
\\[\AgdaEmptyExtraSkip]%
\>[0]\AgdaKeyword{private}\AgdaSpace{}%
\AgdaKeyword{variable}\<%
\\
\>[0][@{}l@{\AgdaIndent{0}}]%
\>[2]\AgdaGeneralizable{n}\AgdaSpace{}%
\AgdaSymbol{:}\AgdaSpace{}%
\AgdaDatatype{ℕ}\<%
\end{code}
\begin{code}%
\>[0]\AgdaFunction{Ex}\AgdaSpace{}%
\AgdaSymbol{:}\AgdaSpace{}%
\AgdaDatatype{Word}\AgdaSpace{}%
\AgdaSymbol{(}\AgdaPostulate{Gen}\AgdaSpace{}%
\AgdaSymbol{(}\AgdaInductiveConstructor{₂₊}\AgdaSpace{}%
\AgdaGeneralizable{n}\AgdaSymbol{))}\<%
\\
\>[0]\AgdaFunction{Ex}\AgdaSpace{}%
\AgdaSymbol{=}\AgdaSpace{}%
\AgdaPostulate{CZ}\AgdaSpace{}%
\AgdaOperator{\AgdaInductiveConstructor{•}}\AgdaSpace{}%
\AgdaPostulate{H}\AgdaSpace{}%
\AgdaOperator{\AgdaFunction{↓}}\AgdaSpace{}%
\AgdaOperator{\AgdaInductiveConstructor{•}}\AgdaSpace{}%
\AgdaPostulate{H}\AgdaSpace{}%
\AgdaOperator{\AgdaPostulate{↑}}\AgdaSpace{}%
\AgdaOperator{\AgdaInductiveConstructor{•}}\AgdaSpace{}%
\AgdaPostulate{CZ}\AgdaSpace{}%
\AgdaOperator{\AgdaInductiveConstructor{•}}\AgdaSpace{}%
\AgdaPostulate{H}\AgdaSpace{}%
\AgdaOperator{\AgdaFunction{↓}}\AgdaSpace{}%
\AgdaOperator{\AgdaInductiveConstructor{•}}\AgdaSpace{}%
\AgdaPostulate{H}\AgdaSpace{}%
\AgdaOperator{\AgdaPostulate{↑}}\AgdaSpace{}%
\AgdaOperator{\AgdaInductiveConstructor{•}}\AgdaSpace{}%
\AgdaPostulate{CZ}\AgdaSpace{}%
\AgdaOperator{\AgdaInductiveConstructor{•}}\AgdaSpace{}%
\AgdaPostulate{H}\AgdaSpace{}%
\AgdaOperator{\AgdaFunction{↓}}\AgdaSpace{}%
\AgdaOperator{\AgdaInductiveConstructor{•}}\AgdaSpace{}%
\AgdaPostulate{H}\AgdaSpace{}%
\AgdaOperator{\AgdaPostulate{↑}}\<%
\end{code}

\end{minipage}\hfill
\begin{minipage}[c]{0.42\textwidth}\centering
\ufig[0.48]{def-Ex-eq}
\end{minipage}

\noindent A rule set is an inductive family indexed by the width, with one constructor per rule. Each rule is stated once, on the bottom wires of every width at which it fits. For example, these are two of the fifteen constructors of our primary rule set \aid{{\_}SRel,{\_}==={\_}}, the order of the swap and the rule moving $H$ across it:

\noindent\begin{minipage}[c]{0.56\textwidth}
%% Generated by make agda from lagda/ex-relation.lagda.tex (agda --latex --only-scope-checking);
%% edit that file, not this one.  The hidden block renders nothing.
\begin{code}[hide]%
\>[0]\AgdaKeyword{module}\AgdaSpace{}%
\AgdaModule{vqupit-short.lagda.ex-relation}\AgdaSpace{}%
\AgdaKeyword{where}\<%
\\
%
\\[\AgdaEmptyExtraSkip]%
\>[0]\AgdaKeyword{open}\AgdaSpace{}%
\AgdaKeyword{import}\AgdaSpace{}%
\AgdaModule{vqupit-short.lagda.Prelude}\<%
\\
\>[0]\AgdaKeyword{open}\AgdaSpace{}%
\AgdaKeyword{import}\AgdaSpace{}%
\AgdaModule{vqupit-short.lagda.ex-circuit}\<%
\\
%
\\[\AgdaEmptyExtraSkip]%
\>[0]\AgdaKeyword{private}\AgdaSpace{}%
\AgdaKeyword{variable}\<%
\\
\>[0][@{}l@{\AgdaIndent{0}}]%
\>[2]\AgdaGeneralizable{n}\AgdaSpace{}%
\AgdaSymbol{:}\AgdaSpace{}%
\AgdaDatatype{ℕ}\<%
\\
%
\\[\AgdaEmptyExtraSkip]%
\>[0]\AgdaKeyword{infix}\AgdaSpace{}%
\AgdaNumber{4}\AgdaSpace{}%
\AgdaOperator{\AgdaDatatype{\AgdaUnderscore{}SRel,\AgdaUnderscore{}===\AgdaUnderscore{}}}\<%
\end{code}
\begin{code}%
\>[0]\AgdaKeyword{data}\AgdaSpace{}%
\AgdaOperator{\AgdaDatatype{\AgdaUnderscore{}SRel,\AgdaUnderscore{}===\AgdaUnderscore{}}}\AgdaSpace{}%
\AgdaSymbol{:}\AgdaSpace{}%
\AgdaSymbol{(}\AgdaBound{n}\AgdaSpace{}%
\AgdaSymbol{:}\AgdaSpace{}%
\AgdaDatatype{ℕ}\AgdaSymbol{)}\AgdaSpace{}%
\AgdaSymbol{→}\AgdaSpace{}%
\AgdaFunction{WRel}\AgdaSpace{}%
\AgdaSymbol{(}\AgdaPostulate{Gen}\AgdaSpace{}%
\AgdaBound{n}\AgdaSymbol{)}\AgdaSpace{}%
\AgdaKeyword{where}\<%
\\
\>[0][@{}l@{\AgdaIndent{0}}]%
\>[2]\AgdaInductiveConstructor{order-Ex}%
\>[14]\AgdaSymbol{:}\AgdaSpace{}%
\AgdaSymbol{(}\AgdaInductiveConstructor{₂₊}\AgdaSpace{}%
\AgdaGeneralizable{n}\AgdaSymbol{)}\AgdaSpace{}%
\AgdaOperator{\AgdaDatatype{SRel,}}%
\>[30]\AgdaFunction{Ex}\AgdaSpace{}%
\AgdaOperator{\AgdaFunction{\textasciicircum{}}}\AgdaSpace{}%
\AgdaNumber{2}\AgdaSpace{}%
\AgdaOperator{\AgdaDatatype{===}}\AgdaSpace{}%
\AgdaInductiveConstructor{ε}\<%
\\
%
\>[2]\AgdaInductiveConstructor{semi-Ex-H↑}%
\>[14]\AgdaSymbol{:}\AgdaSpace{}%
\AgdaSymbol{(}\AgdaInductiveConstructor{₂₊}\AgdaSpace{}%
\AgdaGeneralizable{n}\AgdaSymbol{)}\AgdaSpace{}%
\AgdaOperator{\AgdaDatatype{SRel,}}%
\>[30]\AgdaFunction{Ex}\AgdaSpace{}%
\AgdaOperator{\AgdaInductiveConstructor{•}}\AgdaSpace{}%
\AgdaPostulate{H}\AgdaSpace{}%
\AgdaOperator{\AgdaPostulate{↑}}\AgdaSpace{}%
\AgdaOperator{\AgdaDatatype{===}}\AgdaSpace{}%
\AgdaPostulate{H}\AgdaSpace{}%
\AgdaOperator{\AgdaFunction{↓}}\AgdaSpace{}%
\AgdaOperator{\AgdaInductiveConstructor{•}}\AgdaSpace{}%
\AgdaFunction{Ex}\<%
\end{code}

\end{minipage}\hfill
\begin{minipage}[c]{0.42\textwidth}\centering
\ufig[0.48]{pv0-order-Ex}\\[3pt]
\ufig[0.48]{pv0-semi-Ex-Hup}
\end{minipage}

\noindent Reasoning with the rules is equational reasoning in the congruence $\approx$ they generate, with combinators for symmetry, associativity, and rewriting the left or right factor of a product. For example, conjugating by the swap moves $H$ down a wire:

\noindent\begin{minipage}[c]{0.56\textwidth}
%% Generated by make agda from lagda/ex-reasoning.lagda.tex (agda --latex --only-scope-checking);
%% edit that file, not this one.  The hidden block renders nothing.
\begin{code}[hide]%
\>[0]\AgdaKeyword{module}\AgdaSpace{}%
\AgdaModule{vqupit-short.lagda.ex-reasoning}\AgdaSpace{}%
\AgdaKeyword{where}\<%
\\
%
\\[\AgdaEmptyExtraSkip]%
\>[0]\AgdaKeyword{open}\AgdaSpace{}%
\AgdaKeyword{import}\AgdaSpace{}%
\AgdaModule{vqupit-short.lagda.Prelude}\<%
\\
\>[0]\AgdaKeyword{open}\AgdaSpace{}%
\AgdaKeyword{import}\AgdaSpace{}%
\AgdaModule{vqupit-short.lagda.ex-circuit}\<%
\\
\>[0]\AgdaKeyword{open}\AgdaSpace{}%
\AgdaKeyword{import}\AgdaSpace{}%
\AgdaModule{vqupit-short.lagda.ex-relation}\<%
\\
%
\\[\AgdaEmptyExtraSkip]%
\>[0]\AgdaKeyword{private}\AgdaSpace{}%
\AgdaKeyword{variable}\<%
\\
\>[0][@{}l@{\AgdaIndent{0}}]%
\>[2]\AgdaGeneralizable{n}\AgdaSpace{}%
\AgdaSymbol{:}\AgdaSpace{}%
\AgdaDatatype{ℕ}\<%
\\
%
\\[\AgdaEmptyExtraSkip]%
\>[0]\AgdaKeyword{postulate}\<%
\\
\>[0][@{}l@{\AgdaIndent{0}}]%
\>[2]\AgdaPostulate{axiom}\AgdaSpace{}%
\AgdaSymbol{:}\AgdaSpace{}%
\AgdaSymbol{∀}\AgdaSpace{}%
\AgdaSymbol{\{}\AgdaBound{w}\AgdaSpace{}%
\AgdaBound{v}\AgdaSpace{}%
\AgdaSymbol{:}\AgdaSpace{}%
\AgdaDatatype{Word}\AgdaSpace{}%
\AgdaSymbol{(}\AgdaPostulate{Gen}\AgdaSpace{}%
\AgdaGeneralizable{n}\AgdaSymbol{)\}}\AgdaSpace{}%
\AgdaSymbol{→}\AgdaSpace{}%
\AgdaGeneralizable{n}\AgdaSpace{}%
\AgdaOperator{\AgdaDatatype{SRel,}}\AgdaSpace{}%
\AgdaBound{w}\AgdaSpace{}%
\AgdaOperator{\AgdaDatatype{===}}\AgdaSpace{}%
\AgdaBound{v}\AgdaSpace{}%
\AgdaSymbol{→}\AgdaSpace{}%
\AgdaBound{w}\AgdaSpace{}%
\AgdaOperator{\AgdaPostulate{≈}}\AgdaSpace{}%
\AgdaBound{v}\<%
\end{code}
\begin{code}%
\>[0]\AgdaFunction{conj-Ex-H↑}\AgdaSpace{}%
\AgdaSymbol{:}\AgdaSpace{}%
\AgdaFunction{Ex}\AgdaSpace{}%
\AgdaOperator{\AgdaInductiveConstructor{•}}\AgdaSpace{}%
\AgdaPostulate{H}\AgdaSpace{}%
\AgdaOperator{\AgdaPostulate{↑}}\AgdaSpace{}%
\AgdaOperator{\AgdaInductiveConstructor{•}}\AgdaSpace{}%
\AgdaFunction{Ex}\AgdaSpace{}%
\AgdaOperator{\AgdaPostulate{≈}}\AgdaSpace{}%
\AgdaPostulate{H}\AgdaSpace{}%
\AgdaOperator{\AgdaFunction{↓}}\<%
\\
\>[0]\AgdaFunction{conj-Ex-H↑}\AgdaSpace{}%
\AgdaSymbol{=}\AgdaSpace{}%
\AgdaOperator{\AgdaPostulate{begin}}\<%
\\
\>[0][@{}l@{\AgdaIndent{0}}]%
\>[2]\AgdaFunction{Ex}\AgdaSpace{}%
\AgdaOperator{\AgdaInductiveConstructor{•}}\AgdaSpace{}%
\AgdaPostulate{H}\AgdaSpace{}%
\AgdaOperator{\AgdaPostulate{↑}}\AgdaSpace{}%
\AgdaOperator{\AgdaInductiveConstructor{•}}\AgdaSpace{}%
\AgdaFunction{Ex}%
\>[19]\AgdaOperator{\AgdaPostulate{≈⟨}}\AgdaSpace{}%
\AgdaPostulate{sym}\AgdaSpace{}%
\AgdaPostulate{assoc}\AgdaSpace{}%
\AgdaOperator{\AgdaPostulate{⟩}}\<%
\\
%
\>[2]\AgdaSymbol{(}\AgdaFunction{Ex}\AgdaSpace{}%
\AgdaOperator{\AgdaInductiveConstructor{•}}\AgdaSpace{}%
\AgdaPostulate{H}\AgdaSpace{}%
\AgdaOperator{\AgdaPostulate{↑}}\AgdaSymbol{)}\AgdaSpace{}%
\AgdaOperator{\AgdaInductiveConstructor{•}}\AgdaSpace{}%
\AgdaFunction{Ex}%
\>[19]\AgdaOperator{\AgdaPostulate{≈⟨}}\AgdaSpace{}%
\AgdaOperator{\AgdaPostulate{cleft}}\AgdaSpace{}%
\AgdaPostulate{axiom}\AgdaSpace{}%
\AgdaInductiveConstructor{semi-Ex-H↑}\AgdaSpace{}%
\AgdaOperator{\AgdaPostulate{⟩}}\<%
\\
%
\>[2]\AgdaSymbol{(}\AgdaPostulate{H}\AgdaSpace{}%
\AgdaOperator{\AgdaFunction{↓}}\AgdaSpace{}%
\AgdaOperator{\AgdaInductiveConstructor{•}}\AgdaSpace{}%
\AgdaFunction{Ex}\AgdaSymbol{)}\AgdaSpace{}%
\AgdaOperator{\AgdaInductiveConstructor{•}}\AgdaSpace{}%
\AgdaFunction{Ex}%
\>[19]\AgdaOperator{\AgdaPostulate{≈⟨}}\AgdaSpace{}%
\AgdaPostulate{assoc}\AgdaSpace{}%
\AgdaOperator{\AgdaPostulate{⟩}}\<%
\\
%
\>[2]\AgdaPostulate{H}\AgdaSpace{}%
\AgdaOperator{\AgdaFunction{↓}}\AgdaSpace{}%
\AgdaOperator{\AgdaInductiveConstructor{•}}\AgdaSpace{}%
\AgdaFunction{Ex}\AgdaSpace{}%
\AgdaOperator{\AgdaInductiveConstructor{•}}\AgdaSpace{}%
\AgdaFunction{Ex}%
\>[19]\AgdaOperator{\AgdaPostulate{≈⟨}}\AgdaSpace{}%
\AgdaOperator{\AgdaPostulate{cright}}\AgdaSpace{}%
\AgdaPostulate{axiom}\AgdaSpace{}%
\AgdaInductiveConstructor{order-Ex}\AgdaSpace{}%
\AgdaOperator{\AgdaPostulate{⟩}}\<%
\\
%
\>[2]\AgdaPostulate{H}\AgdaSpace{}%
\AgdaOperator{\AgdaFunction{↓}}\AgdaSpace{}%
\AgdaOperator{\AgdaInductiveConstructor{•}}\AgdaSpace{}%
\AgdaInductiveConstructor{ε}%
\>[19]\AgdaOperator{\AgdaPostulate{≈⟨}}\AgdaSpace{}%
\AgdaPostulate{right-unit}\AgdaSpace{}%
\AgdaOperator{\AgdaPostulate{⟩}}\<%
\\
%
\>[2]\AgdaPostulate{H}\AgdaSpace{}%
\AgdaOperator{\AgdaFunction{↓}}\AgdaSpace{}%
\AgdaOperator{\AgdaPostulate{∎}}\<%
\end{code}

\end{minipage}\hfill
\begin{minipage}[c]{0.42\textwidth}\centering
\ufig[0.48]{pv0-conj-Ex-Hup-steps}
\end{minipage}

\noindent We can then state our main theorems:

\begin{theorem}[\aid{Paper-V1.presentation}]\label{thm:q-main}
For every width $n$, the fifteen rules of \texttt{Paper-V1} (\Cref{sec:q-composing}) present the projective Clifford group:
%% Copied from ../vqupit/latex/vqupit/sections/composing.tex (composing, the clifford-presentation block), the LaTeX
%% that `agda --latex --only-scope-checking' generated for the long version.
\begin{code}%
\>[0]\AgdaFunction{presentation}\AgdaSpace{}%
\AgdaSymbol{:}\AgdaSpace{}%
\AgdaSymbol{∀}\AgdaSpace{}%
\AgdaBound{n}\AgdaSpace{}%
\AgdaSymbol{→}\AgdaSpace{}%
\AgdaSymbol{(}\AgdaOperator{\AgdaPostulate{PapV1.Clifford-Relations.\AgdaUnderscore{}QRel,\AgdaUnderscore{}===\AgdaUnderscore{}}}\AgdaSpace{}%
\AgdaBound{n}\AgdaSymbol{)}\AgdaSpace{}%
\AgdaOperator{\AgdaRecord{IsPresentationOf}}\AgdaSpace{}%
\AgdaSymbol{(}\AgdaPostulate{Pauli⋊Sp}\AgdaSpace{}%
\AgdaBound{n}\AgdaSymbol{)}\<%
\end{code}

\end{theorem}

\noindent These rules are the relations of~\quditcite{} read modulo scalars, with its rule~(C4) in an equivalent form. The relation \aid{{\_}QRel,{\_}==={\_}} closes \aid{{\_}SRel,{\_}==={\_}} under the structural rules (\Cref{sec:form}), and $\aid{Pauli⋊Sp}\;n$ is the semidirect product of $\Zp^{2n}$ by $\Sp$ (\Cref{sec:q-composing}). The record \aid{{\_}IsPresentationOf{\_}} (\Cref{sec:form}) gives a group isomorphism from circuits modulo $\approx$ to $\PPauli \rtimes \Sp$. That this map is well defined is \emph{soundness}, since each rule holds. That it is injective is \emph{completeness}, and that it is surjective says that every projective Clifford operator is represented by a circuit.

We refine the normal form of~\quditcite{} for symplectic operators, which are Clifford operators modulo Paulis, into a doubly inductive definition (\Cref{sec:q-symplectic}), and we prove that normal forms are unique:

\begin{theorem}[\aid{unique-nf}]\label{thm:q-unique}
Normal forms with equal symplectic denotations are equal:
%% Copied from ../vqupit/latex/vqupit/sections/symplectic.tex (symplectic, the unique-nf block), the LaTeX
%% that `agda --latex --only-scope-checking' generated for the long version.
\begin{code}%
\>[0]\AgdaFunction{unique-nf}\AgdaSpace{}%
\AgdaSymbol{:}\AgdaSpace{}%
\AgdaSymbol{∀}\AgdaSpace{}%
\AgdaBound{n}\AgdaSpace{}%
\AgdaSymbol{\{}\AgdaBound{u}\AgdaSpace{}%
\AgdaBound{v}\AgdaSpace{}%
\AgdaSymbol{:}\AgdaSpace{}%
\AgdaFunction{NF}\AgdaSpace{}%
\AgdaBound{n}\AgdaSymbol{\}}\AgdaSpace{}%
\AgdaSymbol{→}\AgdaSpace{}%
\AgdaOperator{\AgdaPostulate{⟦}}\AgdaSpace{}%
\AgdaOperator{\AgdaFunction{[}}\AgdaSpace{}%
\AgdaBound{u}\AgdaSpace{}%
\AgdaOperator{\AgdaFunction{]}}\AgdaSpace{}%
\AgdaOperator{\AgdaPostulate{⟧}}\AgdaSpace{}%
\AgdaOperator{\AgdaPostulate{≈ˢ}}\AgdaSpace{}%
\AgdaOperator{\AgdaPostulate{⟦}}\AgdaSpace{}%
\AgdaOperator{\AgdaFunction{[}}\AgdaSpace{}%
\AgdaBound{v}\AgdaSpace{}%
\AgdaOperator{\AgdaFunction{]}}\AgdaSpace{}%
\AgdaOperator{\AgdaPostulate{⟧}}\AgdaSpace{}%
\AgdaSymbol{→}\AgdaSpace{}%
\AgdaBound{u}\AgdaSpace{}%
\AgdaOperator{\AgdaDatatype{≡}}\AgdaSpace{}%
\AgdaBound{v}\<%
\end{code}

\end{theorem}

\paragraph{The route} We mostly follow the paper. Its normal boxes, its symplectic normal form, and its parametrized \emph{box relations} reappear below. A box relation says that pushing a gate through a normal box yields an updated box and residual \emph{dirty} gates. We also keep the reduction to eighteen relations of its Theorem~4.4, a step the paper already delegated to Agda. The one major divergence is the step that promotes symplectic to projective completeness, its Proposition~4.9. There, the paper corrects each rule $C_l = C_r$ to $C_l = C_r\,;P$ for the unique Pauli $P$ that makes it projectively sound. It then replays the symplectic rewrites, pushing each correction to the end, where it becomes ``a potentially different Pauli''. We did not see how to verify this argument at reasonable cost. Instead, we work in a syntax where Paulis are trivial, since symplectic circuits have only $H$, $S$, and $CZ$ and denote $\Sp$ itself. We then reintroduce Paulis and scalars by construction (\Cref{fig:q-route}). The Pauli bookkeeping of Proposition~4.9 is thus done once, by a theorem about semidirect products that the framework proves for arbitrary groups.

\begin{figure}[tbp]
\centering
% Qudit presentation route. Folded into two columns to keep labels readable.
\begingroup
\definecolor{routeBlue}{HTML}{205886}
\definecolor{routeInk}{HTML}{233447}
\begin{tikzpicture}[
  x=1cm,y=1cm,line cap=round,line join=round,
  route box/.style={draw=routeBlue!60,fill=white,rounded corners=1.5mm,
    line width=.65pt,text=routeInk,align=center,
    font=\sffamily\fontsize{8}{9.5}\selectfont,
    text width=4.10cm,minimum height=1.03cm,inner sep=4pt},
  route label/.style={font=\sffamily\fontsize{7.4}{8.7}\selectfont,
    text=routeInk,align=center,fill=white,inner sep=1.5pt},
  route arrow/.style={-{Stealth[length=1.6mm,width=1.4mm]},
    draw=routeBlue,line width=.85pt}
]
% Symplectic presentations, followed by the Pauli extension.
\node[route box] (sp) at (2.25,0)
  {$\Sp$: full rules (17)\\coset tower, unique NF};
\node[route box] (sim) at (11.25,0)
  {$\Sp$: simplified\\rules (15)};
\node[route box] (pauli) at (2.25,-2)
  {$\PPauli$\\$X^p,\ Z^p,\ XZ=ZX$};
\node[route box] (sd) at (11.25,-2)
  {$\PPauli\rtimes\Sp$\\semidirect presentation\\on $X,Z,H,S,CZ$};

% Generator reduction and the two equivalent projective rule sets.
\node[route box] (v0) at (2.25,-4)
  {Paper-V0 (16 rules)\\multiplier over $R$};
\node[route box] (v1) at (11.25,-4)
  {Simplified-V1\\19 rules over $H,S,CZ$};
\node[route box] (pv1) at (2.25,-6)
  {Paper-V1 (15 rules)\\Theorem~\ref{thm:q-main}\\the paper's Figure~1 mod scalars};
\node[route box] (ex) at (11.25,-6)
  {Exact: Paper-V0 $+\,\omega$\\one correction};

% The scalar extension is a separate branch from Paper-V0.
\node[route box] (fig1) at (2.25,-8)
  {the paper's Figure~1\\16 rules over $-\omega$};
\node[route box] (exm) at (11.25,-8)
  {Exact $+\,{-1}$};

\draw[route arrow] (sp.east) --
  node[route label,above=3pt] {identity isomorphism} (sim.west);
\draw[route arrow] (sim.south) -- (sd.north);
\node[route label,anchor=east,text width=2.7cm] at (11.00,-.96)
  {semidirect\\theorem};
\draw[route arrow] (pauli.east) --
  node[route label,above=3pt] {semidirect theorem} (sd.west);
\draw[route arrow] (sd.south) -- (v1.north);
\node[route label,anchor=east,text width=2.7cm] at (11.00,-3.04)
  {generator reduction:\\$X,Z\mapsto$ words};
\draw[route arrow] (v1.west) --
  node[route label,above=3pt] {identity isomorphism} (v0.east);
\draw[route arrow] (v0.south) -- (pv1.north);
\node[route label,anchor=east,text width=2.05cm] at (2.00,-4.96)
  {identity\\isomorphism};
\draw[route arrow] (v0.south east) -- (ex.north west);
\node[route label,text width=3.1cm] at (6.75,-4.40)
  {central extension theorem};
\draw[route arrow] (ex.south) -- (exm.north);
\node[route label,anchor=east,text width=2.7cm] at (11.00,-7)
  {direct product\\theorem};
\draw[route arrow] (exm.west) --
  node[route label,above=3pt] {isomorphism:\\$-\omega\mapsto-1\cdot\omega$} (fig1.east);
\end{tikzpicture}
\endgroup

\caption{The chain of presentations used to formalize the qudit Clifford result of~\quditcite{}. Every arrow is a verified presentation theorem or isomorphism of presentations. The symplectic stages present $\Sp$, and the projective chain from the semidirect presentation to Paper-V1 presents $\PPauli\rtimes\Sp$. ``Exact'' is a central extension by $\langle\omega\rangle$, and the final two stages also include the direct factor $\langle-1\rangle$.}
\Description{A two-column flow connects the full and simplified symplectic presentations, then combines the latter with a Pauli presentation by a semidirect-product theorem. Generator reduction yields Simplified-V1, followed by Paper-V0. One branch reaches Paper-V1 modulo scalars. The other adds scalar phases by a central extension, then adjoins minus one by a direct product and reaches the paper's Figure 1 by an isomorphism.}
\label{fig:q-route}
\end{figure}

\subsection{The Symplectic Factor}\label{sec:q-symplectic}

We model $\Sp$ as the group of $\Zp$-linear bijections of the phase space $\aid{Pauli}\;n = \aid{Vec}\;(\Zp \times \Zp)\;n$ that preserve the symplectic form. An element of the phase space lists, for each wire, the exponents of $X^{a}Z^{b}$ up to phase. Against this semantics, completeness rests on the normal form, which we refine into a doubly inductive datatype:

\noindent\begin{minipage}[c]{0.4\textwidth}
%% Copied from ../vqupit/latex/vqupit/sections/symplectic.tex (symplectic, block 2), the LaTeX
%% that `agda --latex --only-scope-checking' generated for the long version.
\begin{code}%
\>[0]\AgdaFunction{M}\AgdaSpace{}%
\AgdaFunction{L}\AgdaSpace{}%
\AgdaFunction{ML}\AgdaSpace{}%
\AgdaSymbol{:}\AgdaSpace{}%
\AgdaDatatype{ℕ}\AgdaSpace{}%
\AgdaSymbol{→}\AgdaSpace{}%
\AgdaPrimitive{Set}\<%
\\
\>[0]\AgdaFunction{M}\AgdaSpace{}%
\AgdaNumber{0}\AgdaSpace{}%
\AgdaSymbol{=}\AgdaSpace{}%
\AgdaRecord{⊤}\<%
\\
\>[0]\AgdaFunction{M}\AgdaSpace{}%
\AgdaSymbol{(}\AgdaInductiveConstructor{₁₊}\AgdaSpace{}%
\AgdaBound{n}\AgdaSymbol{)}\AgdaSpace{}%
\AgdaSymbol{=}\AgdaSpace{}%
\AgdaDatatype{Vec}\AgdaSpace{}%
\AgdaPostulate{D}\AgdaSpace{}%
\AgdaBound{n}\AgdaSpace{}%
\AgdaOperator{\AgdaFunction{×}}\AgdaSpace{}%
\AgdaPostulate{E}\<%
\\
\>[0]\AgdaFunction{L}\AgdaSpace{}%
\AgdaNumber{0}\AgdaSpace{}%
\AgdaSymbol{=}\AgdaSpace{}%
\AgdaRecord{⊤}\<%
\\
\>[0]\AgdaFunction{L}\AgdaSpace{}%
\AgdaSymbol{(}\AgdaInductiveConstructor{₁₊}\AgdaSpace{}%
\AgdaBound{n}\AgdaSymbol{)}\AgdaSpace{}%
\AgdaSymbol{=}\AgdaSpace{}%
\AgdaDatatype{Vec}\AgdaSpace{}%
\AgdaPostulate{B}\AgdaSpace{}%
\AgdaBound{n}\AgdaSpace{}%
\AgdaOperator{\AgdaFunction{×}}\AgdaSpace{}%
\AgdaPostulate{A}\<%
\\
\>[0]\AgdaFunction{ML}\AgdaSpace{}%
\AgdaNumber{0}\AgdaSpace{}%
\AgdaSymbol{=}\AgdaSpace{}%
\AgdaRecord{⊤}\<%
\\
\>[0]\AgdaFunction{ML}\AgdaSpace{}%
\AgdaNumber{1}\AgdaSpace{}%
\AgdaSymbol{=}\AgdaSpace{}%
\AgdaFunction{M}\AgdaSpace{}%
\AgdaNumber{1}\AgdaSpace{}%
\AgdaOperator{\AgdaFunction{×}}\AgdaSpace{}%
\AgdaFunction{L}\AgdaSpace{}%
\AgdaNumber{1}\<%
\\
\>[0]\AgdaFunction{ML}\AgdaSpace{}%
\AgdaBound{m}\AgdaSymbol{@(}\AgdaInductiveConstructor{₂₊}\AgdaSpace{}%
\AgdaBound{n}\AgdaSymbol{)}%
\>[350I]\AgdaSymbol{=}\AgdaSpace{}%
\AgdaFunction{M}\AgdaSpace{}%
\AgdaBound{m}\AgdaSpace{}%
\AgdaOperator{\AgdaFunction{×}}\AgdaSpace{}%
\AgdaFunction{L}\AgdaSpace{}%
\AgdaBound{m}\<%
\\
\>[.][@{}l@{}]\<[350I]%
\>[12]\AgdaOperator{\AgdaDatatype{⊎}}\AgdaSpace{}%
\AgdaPostulate{D}\AgdaSpace{}%
\AgdaOperator{\AgdaFunction{×}}\AgdaSpace{}%
\AgdaFunction{ML}\AgdaSpace{}%
\AgdaSymbol{(}\AgdaInductiveConstructor{₁₊}\AgdaSpace{}%
\AgdaBound{n}\AgdaSymbol{)}\<%
\end{code}

\end{minipage}\hfill
\begin{minipage}[c]{0.56\textwidth}\centering
  \begin{tabular}{@{}r@{\,}c@{\,}l@{\quad}r@{\,}c@{\,}l@{}}
    \ufig[0.5]{m4-box} & \scalebox{0.7}{$=$} & \ufig[0.5]{m4-atoms} &
    \ufig[0.5]{ml4-box} & \scalebox{0.7}{$=$} & \ufig[0.5]{ml-tower-a-r}\\[1ex]
    \ufig[0.5]{l4-box} & \scalebox{0.7}{$=$} & \ufig[0.5]{l4-atoms} &
    \ufig[0.5]{ml4-box} & \scalebox{0.7}{$=$} & \ufig[0.5]{ml-tower-b-r}\\
  \end{tabular}
\end{minipage}

\medskip\noindent
An $M$-row is a column of $D$ boxes ending in an $E$ box, the paper's $X$-normal layer. An $L$-row is a column of $B$ boxes ending in an $A$ box, its $Z$-normal layer, with the $A$ box at the bottom. An $\ML$ box at width $m \ge 2$ is either a full pair of rows or a bottom $D$ box followed by an $\ML$ box one wire up, so it is a spine of $D$ boxes capped by a pair of rows. This is the first induction, on the width of a single coset representative. The second induction is the staircase:

\noindent\begin{minipage}[c]{0.54\textwidth}
%% Copied from ../vqupit/latex/vqupit/sections/symplectic.tex (symplectic, block 3), the LaTeX
%% that `agda --latex --only-scope-checking' generated for the long version.
\begin{code}%
\>[0]\AgdaFunction{NF}\AgdaSpace{}%
\AgdaSymbol{:}\AgdaSpace{}%
\AgdaDatatype{ℕ}\AgdaSpace{}%
\AgdaSymbol{→}\AgdaSpace{}%
\AgdaPrimitive{Set}\<%
\\
\>[0]\AgdaFunction{NF}\AgdaSpace{}%
\AgdaNumber{0}\AgdaSpace{}%
\AgdaSymbol{=}\AgdaSpace{}%
\AgdaRecord{⊤}\<%
\\
\>[0]\AgdaFunction{NF}\AgdaSpace{}%
\AgdaSymbol{(}\AgdaInductiveConstructor{₁₊}\AgdaSpace{}%
\AgdaBound{n}\AgdaSymbol{)}\AgdaSpace{}%
\AgdaSymbol{=}\AgdaSpace{}%
\AgdaFunction{NF}\AgdaSpace{}%
\AgdaBound{n}\AgdaSpace{}%
\AgdaOperator{\AgdaFunction{×}}\AgdaSpace{}%
\AgdaFunction{ML}\AgdaSpace{}%
\AgdaSymbol{(}\AgdaInductiveConstructor{₁₊}\AgdaSpace{}%
\AgdaBound{n}\AgdaSymbol{)}\<%
\\
%
\\[\AgdaEmptyExtraSkip]%
\>[0]\AgdaOperator{\AgdaFunction{[\AgdaUnderscore{}]}}\AgdaSpace{}%
\AgdaSymbol{:}\AgdaSpace{}%
\AgdaSymbol{∀}\AgdaSpace{}%
\AgdaSymbol{\{}\AgdaBound{n}\AgdaSymbol{\}}\AgdaSpace{}%
\AgdaSymbol{→}\AgdaSpace{}%
\AgdaFunction{NF}\AgdaSpace{}%
\AgdaBound{n}\AgdaSpace{}%
\AgdaSymbol{→}\AgdaSpace{}%
\AgdaDatatype{Word}\AgdaSpace{}%
\AgdaSymbol{(}\AgdaDatatype{Gen}\AgdaSpace{}%
\AgdaBound{n}\AgdaSymbol{)}\<%
\\
\>[0]\AgdaOperator{\AgdaFunction{[\AgdaUnderscore{}]}}\AgdaSpace{}%
\AgdaSymbol{\{}\AgdaNumber{0}\AgdaSymbol{\}}%
\>[12]\AgdaInductiveConstructor{tt}%
\>[23]\AgdaSymbol{=}\AgdaSpace{}%
\AgdaInductiveConstructor{ε}\<%
\\
\>[0]\AgdaOperator{\AgdaFunction{[\AgdaUnderscore{}]}}\AgdaSpace{}%
\AgdaSymbol{\{}\AgdaInductiveConstructor{₁₊}\AgdaSpace{}%
\AgdaBound{n}\AgdaSymbol{\}}%
\>[12]\AgdaSymbol{(}\AgdaBound{nf}\AgdaSpace{}%
\AgdaOperator{\AgdaInductiveConstructor{,}}\AgdaSpace{}%
\AgdaBound{ml}\AgdaSymbol{)}%
\>[23]\AgdaSymbol{=}\AgdaSpace{}%
\AgdaOperator{\AgdaFunction{[}}\AgdaSpace{}%
\AgdaBound{nf}\AgdaSpace{}%
\AgdaOperator{\AgdaFunction{]}}\AgdaSpace{}%
\AgdaOperator{\AgdaPostulate{↑}}\AgdaSpace{}%
\AgdaOperator{\AgdaInductiveConstructor{•}}\AgdaSpace{}%
\AgdaOperator{\AgdaPostulate{[}}\AgdaSpace{}%
\AgdaBound{ml}\AgdaSpace{}%
\AgdaOperator{\AgdaPostulate{]ᵐˡ}}\<%
\end{code}

\end{minipage}\hfill
\begin{minipage}[c]{0.42\textwidth}\centering
  \ufig[0.7]{nf4}
\end{minipage}

\paragraph{The table} The coset table is the paper's ``push a gate through a normal box'', written as a structurally recursive function \aid{ract} with its soundness law \aid{ract-sound} (about 5{,}600 lines together):

%% Generated by make agda from lagda/ex-table.lagda.tex (agda --latex --only-scope-checking);
%% edit that file, not this one.  The hidden block renders nothing.
\begin{code}[hide]%
\>[0]\AgdaKeyword{module}\AgdaSpace{}%
\AgdaModule{vqupit-short.lagda.ex-table}\AgdaSpace{}%
\AgdaKeyword{where}\<%
\\
%
\\[\AgdaEmptyExtraSkip]%
\>[0]\AgdaKeyword{open}\AgdaSpace{}%
\AgdaKeyword{import}\AgdaSpace{}%
\AgdaModule{Data.Product}\AgdaSpace{}%
\AgdaKeyword{using}\AgdaSpace{}%
\AgdaSymbol{(}\AgdaOperator{\AgdaFunction{\AgdaUnderscore{}×\AgdaUnderscore{}}}\AgdaSymbol{;}\AgdaSpace{}%
\AgdaOperator{\AgdaInductiveConstructor{\AgdaUnderscore{},\AgdaUnderscore{}}}\AgdaSymbol{)}\<%
\\
\>[0]\AgdaKeyword{open}\AgdaSpace{}%
\AgdaKeyword{import}\AgdaSpace{}%
\AgdaModule{Word.Base}\AgdaSpace{}%
\AgdaKeyword{using}\AgdaSpace{}%
\AgdaSymbol{(}\AgdaOperator{\AgdaInductiveConstructor{[\AgdaUnderscore{}]ʷ}}\AgdaSymbol{)}\<%
\\
\>[0]\AgdaKeyword{open}\AgdaSpace{}%
\AgdaKeyword{import}\AgdaSpace{}%
\AgdaModule{vqupit-short.lagda.Prelude}\<%
\\
%
\\[\AgdaEmptyExtraSkip]%
\>[0]\AgdaKeyword{private}\AgdaSpace{}%
\AgdaKeyword{variable}\<%
\\
\>[0][@{}l@{\AgdaIndent{0}}]%
\>[2]\AgdaGeneralizable{n}\AgdaSpace{}%
\AgdaSymbol{:}\AgdaSpace{}%
\AgdaDatatype{ℕ}\<%
\\
%
\\[\AgdaEmptyExtraSkip]%
\>[0]\AgdaFunction{Circuit}\AgdaSpace{}%
\AgdaSymbol{:}\AgdaSpace{}%
\AgdaDatatype{ℕ}\AgdaSpace{}%
\AgdaSymbol{→}\AgdaSpace{}%
\AgdaPrimitive{Set}\<%
\\
\>[0]\AgdaFunction{Circuit}\AgdaSpace{}%
\AgdaBound{n}\AgdaSpace{}%
\AgdaSymbol{=}\AgdaSpace{}%
\AgdaDatatype{Word}\AgdaSpace{}%
\AgdaSymbol{(}\AgdaPostulate{Gen}\AgdaSpace{}%
\AgdaBound{n}\AgdaSymbol{)}\<%
\\
%
\\[\AgdaEmptyExtraSkip]%
\>[0]\AgdaKeyword{postulate}\<%
\\
\>[0][@{}l@{\AgdaIndent{0}}]%
\>[2]\AgdaPostulate{ML}\AgdaSpace{}%
\AgdaSymbol{:}\AgdaSpace{}%
\AgdaDatatype{ℕ}\AgdaSpace{}%
\AgdaSymbol{→}\AgdaSpace{}%
\AgdaPrimitive{Set}\<%
\\
%
\>[2]\AgdaOperator{\AgdaPostulate{[\AgdaUnderscore{}]ᶜ}}\AgdaSpace{}%
\AgdaSymbol{:}\AgdaSpace{}%
\AgdaPostulate{ML}\AgdaSpace{}%
\AgdaSymbol{(}\AgdaInductiveConstructor{₁₊}\AgdaSpace{}%
\AgdaGeneralizable{n}\AgdaSymbol{)}\AgdaSpace{}%
\AgdaSymbol{→}\AgdaSpace{}%
\AgdaFunction{Circuit}\AgdaSpace{}%
\AgdaSymbol{(}\AgdaInductiveConstructor{₁₊}\AgdaSpace{}%
\AgdaGeneralizable{n}\AgdaSymbol{)}\<%
\end{code}
\begin{code}%
%
\>[2]\AgdaPostulate{ract}%
\>[14]\AgdaSymbol{:}\AgdaSpace{}%
\AgdaPostulate{ML}\AgdaSpace{}%
\AgdaSymbol{(}\AgdaInductiveConstructor{₁₊}\AgdaSpace{}%
\AgdaGeneralizable{n}\AgdaSymbol{)}\AgdaSpace{}%
\AgdaSymbol{→}\AgdaSpace{}%
\AgdaPostulate{Gen}\AgdaSpace{}%
\AgdaSymbol{(}\AgdaInductiveConstructor{₁₊}\AgdaSpace{}%
\AgdaGeneralizable{n}\AgdaSymbol{)}\AgdaSpace{}%
\AgdaSymbol{→}\AgdaSpace{}%
\AgdaFunction{Circuit}\AgdaSpace{}%
\AgdaGeneralizable{n}\AgdaSpace{}%
\AgdaOperator{\AgdaFunction{×}}\AgdaSpace{}%
\AgdaPostulate{ML}\AgdaSpace{}%
\AgdaSymbol{(}\AgdaInductiveConstructor{₁₊}\AgdaSpace{}%
\AgdaGeneralizable{n}\AgdaSymbol{)}\<%
\\
%
\>[2]\AgdaPostulate{ract-sound}%
\>[14]\AgdaSymbol{:}\AgdaSpace{}%
\AgdaSymbol{∀}\AgdaSpace{}%
\AgdaSymbol{(}\AgdaBound{c}\AgdaSpace{}%
\AgdaSymbol{:}\AgdaSpace{}%
\AgdaPostulate{ML}\AgdaSpace{}%
\AgdaSymbol{(}\AgdaInductiveConstructor{₁₊}\AgdaSpace{}%
\AgdaGeneralizable{n}\AgdaSymbol{))}\AgdaSpace{}%
\AgdaBound{g}\AgdaSpace{}%
\AgdaSymbol{→}\AgdaSpace{}%
\AgdaKeyword{let}\AgdaSpace{}%
\AgdaSymbol{(}\AgdaBound{b'}\AgdaSpace{}%
\AgdaOperator{\AgdaInductiveConstructor{,}}\AgdaSpace{}%
\AgdaBound{c'}\AgdaSymbol{)}\AgdaSpace{}%
\AgdaSymbol{=}\AgdaSpace{}%
\AgdaPostulate{ract}\AgdaSpace{}%
\AgdaBound{c}\AgdaSpace{}%
\AgdaBound{g}\AgdaSpace{}%
\AgdaKeyword{in}\AgdaSpace{}%
\AgdaOperator{\AgdaPostulate{[}}\AgdaSpace{}%
\AgdaBound{c}\AgdaSpace{}%
\AgdaOperator{\AgdaPostulate{]ᶜ}}\AgdaSpace{}%
\AgdaOperator{\AgdaInductiveConstructor{•}}\AgdaSpace{}%
\AgdaOperator{\AgdaInductiveConstructor{[}}\AgdaSpace{}%
\AgdaBound{g}\AgdaSpace{}%
\AgdaOperator{\AgdaInductiveConstructor{]ʷ}}\AgdaSpace{}%
\AgdaOperator{\AgdaPostulate{≈}}\AgdaSpace{}%
\AgdaBound{b'}\AgdaSpace{}%
\AgdaOperator{\AgdaPostulate{↑}}\AgdaSpace{}%
\AgdaOperator{\AgdaInductiveConstructor{•}}\AgdaSpace{}%
\AgdaOperator{\AgdaPostulate{[}}\AgdaSpace{}%
\AgdaBound{c'}\AgdaSpace{}%
\AgdaOperator{\AgdaPostulate{]ᶜ}}\<%
\end{code}


\noindent Each clause of \aid{ract-sound} is derived from the fifteen simplified symplectic relations. For instance, pushing $S$ through an $E$ box is one step on powers of $S$, namely $[\,b\,]^{\mathsf e} \bullet S = S^{-b} \bullet S \approx S^{-(b-1)} = [\,b-1\,]^{\mathsf e}$. A circuit is rewritten to its normal form by repeating this step gate by gate, so every rewrite is traced to a relation.

\paragraph{Uniqueness} The symplectic group acts on $\aid{Vec}\;\mathrm{Pauli}_1\;n$, a $2n$-dimensional symplectic $\Zp$-vector space in which $\mathrm{Pauli}_1$ is the two-dimensional symplectic subspace of a single wire. The elements that fix the bottom wire's subspace pointwise form a copy of $\mathrm{Sp}(2(n{-}1), \Zp)$ acting above the bottom wire, and the $\ML$ boxes represent its cosets. Every $f \in \Sp$ factors as $f = f' \circ c_n$, with $c_n \in \ML_n$ and $f'$ acting above the bottom wire, and $c_n$ is unique because the image of the bottom wire determines the box. Recursing on $f'$ writes $f$ uniquely as $c_1 \circ c_2 \circ \cdots \circ c_n$ with $c_k \in \ML_k$, which is \Cref{thm:q-unique}.

\paragraph{Surjectivity} Every symplectic transformation is the action of some circuit. We prove this constructively by clearing the bottom wire and recursing on the rest. This is the paper's Lemma~3.7, with the coset box computed rather than asserted:

{\renewcommand{\AgdaCodeStyle}{\scriptsize}%% Copied from ../vqupit/latex/vqupit/sections/symplectic.tex (symplectic, the Theorem-LM block), the LaTeX
%% that `agda --latex --only-scope-checking' generated for the long version.
\begin{code}%
\>[0]\AgdaFunction{Theorem-LM}\AgdaSpace{}%
\AgdaSymbol{:}\AgdaSpace{}%
\AgdaSymbol{∀}\AgdaSpace{}%
\AgdaSymbol{\{}\AgdaBound{n}\AgdaSymbol{\}}\AgdaSpace{}%
\AgdaSymbol{(}\AgdaBound{p}\AgdaSpace{}%
\AgdaBound{q}\AgdaSpace{}%
\AgdaSymbol{:}\AgdaSpace{}%
\AgdaPostulate{Pauli}\AgdaSpace{}%
\AgdaBound{n}\AgdaSymbol{)}\AgdaSpace{}%
\AgdaSymbol{→}\AgdaSpace{}%
\AgdaPostulate{sform}\AgdaSpace{}%
\AgdaBound{p}\AgdaSpace{}%
\AgdaBound{q}\AgdaSpace{}%
\AgdaOperator{\AgdaDatatype{≡}}\AgdaSpace{}%
\AgdaInductiveConstructor{₁}\AgdaSpace{}%
\AgdaSymbol{→}\AgdaSpace{}%
\AgdaFunction{∃}\AgdaSpace{}%
\AgdaSymbol{λ}\AgdaSpace{}%
\AgdaSymbol{(}\AgdaBound{lm}\AgdaSpace{}%
\AgdaSymbol{:}\AgdaSpace{}%
\AgdaFunction{ML}\AgdaSpace{}%
\AgdaBound{n}\AgdaSymbol{)}\AgdaSpace{}%
\AgdaSymbol{→}\AgdaSpace{}%
\AgdaPostulate{act}\AgdaSpace{}%
\AgdaOperator{\AgdaPostulate{[}}\AgdaSpace{}%
\AgdaBound{lm}\AgdaSpace{}%
\AgdaOperator{\AgdaPostulate{]ᵐˡ}}\AgdaSpace{}%
\AgdaBound{p}\AgdaSpace{}%
\AgdaOperator{\AgdaDatatype{≡}}\AgdaSpace{}%
\AgdaPostulate{pZ₀}\AgdaSpace{}%
\AgdaOperator{\AgdaFunction{×}}\AgdaSpace{}%
\AgdaPostulate{act}\AgdaSpace{}%
\AgdaOperator{\AgdaPostulate{[}}\AgdaSpace{}%
\AgdaBound{lm}\AgdaSpace{}%
\AgdaOperator{\AgdaPostulate{]ᵐˡ}}\AgdaSpace{}%
\AgdaBound{q}\AgdaSpace{}%
\AgdaOperator{\AgdaDatatype{≡}}\AgdaSpace{}%
\AgdaPostulate{pX₀}\<%
\end{code}
}

\paragraph{Completeness} Circuits with the same denotation rewrite through the table to their normal forms. These normal forms denote the same map by soundness, so they are equal by \Cref{thm:q-unique}, and the circuits are equivalent. This argument is the framework's \aid{by-normalization} (\Cref{fig:normalform}). With surjectivity, it shows that the simplified relations, which their own module also calls \aid{{\_}SRel,{\_}==={\_}}, present the symplectic group:

%% Copied from ../vqupit/latex/vqupit/sections/symplectic.tex (symplectic, the simplified-presentation block), the LaTeX
%% that `agda --latex --only-scope-checking' generated for the long version.
\begin{code}%
\>[0]\AgdaFunction{presentation}\AgdaSpace{}%
\AgdaSymbol{:}\AgdaSpace{}%
\AgdaSymbol{∀}\AgdaSpace{}%
\AgdaBound{n}\AgdaSpace{}%
\AgdaSymbol{→}\AgdaSpace{}%
\AgdaSymbol{(}\AgdaBound{n}\AgdaSpace{}%
\AgdaOperator{\AgdaPostulate{SRel,\AgdaUnderscore{}===\AgdaUnderscore{}}}\AgdaSymbol{)}\AgdaSpace{}%
\AgdaOperator{\AgdaRecord{IsPresentationOf}}\AgdaSpace{}%
\AgdaSymbol{(}\AgdaPostulate{Sp-group}\AgdaSpace{}%
\AgdaBound{n}\AgdaSymbol{)}\<%
\end{code}


\subsection{Composing: Paulis, the Semidirect Product, and Scalars}\label{sec:q-composing}

$\PPauli \cong (\Zp \times \Zp)^n$ is presented over $X$ and $Z$ on each wire by $X^{p} = Z^{p} = \varepsilon$ and $ZX = XZ$. Its semantic group is the phase space that $\Sp$ acts on, and its normal form is $X^{a}Z^{b}$ on each wire. Combining this with the symplectic presentation of \Cref{sec:q-symplectic} is an instance of the framework's theorem for semidirect products. The resulting presentation has the relations of both factors together with the conjugation relations $h\,n\,h^{-1} = \mathit{conj}\,h\,n$, where the action \aid{conj} spells out the paper's Lemmas~2.24--2.25. The theorem asks that the action respect the relations of each factor, and we prove this semantically. The conjugates of congruent words have equal denotations by soundness, so they are congruent by Pauli \emph{completeness}. For symplectic rules this is \aid{respects-Δ}, and it corresponds to the paper's ``Pauli complete'' rules (its Definition~4.7). Since every premise is a theorem, the semidirect presentation of \aid{Pauli⋊Sp} is unconditional.

\paragraph{Reaching the paper's rules} The semidirect presentation has generators $X$, $Z$, $H$, $S$, and $CZ$ and $3 + 15 + 14$ relations (Pauli, symplectic, and conjugation), whereas the paper's rule sets have only $H$, $S$, and $CZ$. Three isomorphisms of presentations close this gap (\Cref{fig:q-route}), through \texttt{Simplified-V1} (nineteen rules), \texttt{Paper-V0} (sixteen), and \texttt{Paper-V1} (the fifteen of \Cref{thm:q-main}). Along the way, $X$ and $Z$ become words, and the Pauli and conjugation rules become theorems.

The first isomorphism, from the semidirect presentation to \texttt{Simplified-V1}, is between different alphabets. Its map $f$ sends $S$ to $R = S \bullet Z^{1/2}$, and its inverse $h$ sends $S$ to $Z^{-1/2} \bullet S$. In particular, $f$ induces the embedding of the symplectic group into the projective Clifford group. The nineteen rules of \texttt{Simplified-V1} are the fifteen simplified symplectic rules together with $(SH)^3 = \varepsilon$, $H^2SH^2S = SH^2SH^2$, and the two rules pushing $X$ through $CZ$. From these, \aid{f-well-defined} derives the images under $f$ of the three Pauli rules and the fourteen conjugation rules.

The second isomorphism, from \texttt{Simplified-V1} to \texttt{Paper-V0}, is the identity on circuits. The two rule sets share six rules word for word and three more up to the spelling of the multiplier, and each derives the rest of the other. The larger direction is \texttt{Paper-V0}'s module \texttt{Lemmas}, of about 4{,}750 lines.

The third isomorphism, from \texttt{Paper-V0} to \texttt{Paper-V1}, changes how the multiplier $M_x\colon \ket{j} \mapsto \ket{xj}$ is spelled as a word. \texttt{Paper-V0} spells it over $R$, as $M_x = R^{x^{-1}}HR^{x}HR^{x^{-1}}H$. The paper's~(T1) spells it over $S$, up to a phase, as $M_x = \aid{SHS'}\;x\;x^{-1}$, where $\aid{SHS'}\;a\;b = Z^{(b-1)/2}X^{(1-a)/2}S^{b}HS^{a}HS^{b}H$. Over $S$, the rule \aid{semi-MR} reads $M_gS^{g^2}Z^{(g^2-g)/2} = SM_g$. \texttt{Paper-V1} restates the four multiplier rules over \aid{SHS'} and drops one rule. Since $M_1$ is now the word $(SH)^3$, \aid{M-power} at $k = 0$ already gives $(SH)^3 \approx \varepsilon$. Over $R$, $M_1$ is the different word $(RH)^3$, so that axiom had to stay. The bridge $R^{a}HR^{b}HR^{a}H \approx \aid{SHS'}\;b\;a$ for $ab = 1$ (\texttt{Shared.PauliBase}) carries the multiplier rules across, which gives \Cref{thm:q-main}. The paper's own rule~(C4) states the conjugation of $S$ by $M_g$ in the other direction. Each of (C4) and \aid{semi-MR} derives the other (\aid{SemiMS}), so Figure~1 of~\quditcite{} read modulo scalars is isomorphic to \texttt{Paper-V1}.

\paragraph{Scalars} The theorem \aid{extension-presentation} extends a projective presentation by a scalar. We apply it to \texttt{Paper-V0}, which has \aid{order-SH} as an axiom. We add $\omega$ with $\omega^{p} = \varepsilon$, let $\omega$ commute with every gate, and \emph{correct} each rule by its scalar. The correction is $\varepsilon$ except on \aid{order-SH}, the one rule that does not hold exactly, since $(SH)^3 = \omega^{-1/8}$. The exponent $-1/8$ is computed in $\Zp$, where it equals $(p^2-1)/8$. The multiplier words now denote $\ket{j} \mapsto \ket{xj}$ only up to a scalar.

The semantic extension is the Weyl cocycle on the natural semidirect product, $c((P,A),(Q,B)) = \tfrac12\,\mathrm{sform}(P,\ A \cdot Q)$. Here $A \cdot Q$ is the symplectic part acting on the Pauli vector, and $\mathrm{sform}$ is the symplectic form of the phase space, the exponent in $PQ = \omega^{\mathrm{sform}(P,Q)}\,QP$. The cocycle identity follows from the bilinearity and the symplectic invariance of \aid{sform}. The cocycle of the presented group is the pullback of $c$ along the presentation by \texttt{Paper-V0}, so the two agree by construction. The identification with unitary operators uses Weyl multiplication~\cite[Eq.~(4), Theorem~3]{gross2006hudson} and a \emph{linear} Weil representation with exact covariance~\cite[Theorem~1.3.2]{gurevich2007geometric}. The operators $U(e,v,A)=\omega^e W(v)\rho(A)$ satisfy
\[
 U(e,v,A)\,U(f,w,B)
 =U\bigl(e+f+\tfrac12\mathrm{sform}(v,Aw),\ v+Aw,\ AB\bigr).
\]
This is the Heisenberg--Weil representation~\cite[\S 1.3]{gurevich2010bounds}. Conjugation on Weyl operators determines $A$ and $v$, so the representation is faithful and its scalars are exactly $\langle\omega I\rangle$ (at width zero, $e\mapsto\omega^e$). Adjoining $-I$ gives its direct product with $\mathbb Z_2$, whose scalars are $\langle-\omega I\rangle$. These finite phase choices differ from the full unitary normalizer, which contains $U(1)$. With no hypotheses, we obtain:

{\renewcommand{\AgdaCodeStyle}{\scriptsize}%% Copied from ../vqupit/latex/vqupit/sections/composing.tex (composing, the presentation-exact block), the LaTeX
%% that `agda --latex --only-scope-checking' generated for the long version.
\begin{code}%
\>[0]\AgdaFunction{presentation-exact}\AgdaSpace{}%
\AgdaSymbol{:}\AgdaSpace{}%
\AgdaSymbol{∀}\AgdaSpace{}%
\AgdaBound{n}\AgdaSpace{}%
\AgdaSymbol{→}\AgdaSpace{}%
\AgdaSymbol{(}\AgdaBound{n}\AgdaSpace{}%
\AgdaOperator{\AgdaPostulate{Exact,\AgdaUnderscore{}===\AgdaUnderscore{}}}\AgdaSymbol{)}\AgdaSpace{}%
\AgdaOperator{\AgdaRecord{IsPresentationOf}}\AgdaSpace{}%
\AgdaSymbol{(}\AgdaPostulate{PE.Presented-group}\AgdaSpace{}%
\AgdaBound{n}\AgdaSpace{}%
\AgdaSymbol{(}\AgdaPostulate{exact-generator-data}\AgdaSpace{}%
\AgdaBound{n}\AgdaSymbol{))}\<%
\end{code}
}

\noindent Agda also shows that the trivial cocycle cannot realize the correction of \aid{order-SH} (\aid{split-forces}).

\paragraph{The rules as printed} Figure~1 of~\quditcite{} is stated over $-\omega$. Its multipliers carry Legendre symbols and its swap carries the sign $(-1)^{(p-1)/2}$, so it needs the $-1$ that the exact group lacks, since the scalars of that group have odd order. Adjoining $-1$ as a central involution presents the exact group times $\mathbb Z_2$, which is \aid{Clifford±}. Figure~1 is isomorphic to this rule set by $-\omega \mapsto -1\cdot\omega$, where the scalar of each exact rule is determined by reading its lift back. This is the paper's Theorem~4.10, stated against our group:

%% Copied from ../vqupit/latex/vqupit/sections/composing.tex (composing, the figure1-presentation block), the LaTeX
%% that `agda --latex --only-scope-checking' generated for the long version.
\begin{code}%
\>[0]\AgdaFunction{presentation-Figure1}\AgdaSpace{}%
\AgdaSymbol{:}\AgdaSpace{}%
\AgdaSymbol{∀}\AgdaSpace{}%
\AgdaBound{n}\AgdaSpace{}%
\AgdaSymbol{→}\AgdaSpace{}%
\AgdaSymbol{(}\AgdaBound{n}\AgdaSpace{}%
\AgdaOperator{\AgdaPostulate{F,\AgdaUnderscore{}===\AgdaUnderscore{}}}\AgdaSymbol{)}\AgdaSpace{}%
\AgdaOperator{\AgdaRecord{IsPresentationOf}}\AgdaSpace{}%
\AgdaSymbol{(}\AgdaPostulate{Clifford±}\AgdaSpace{}%
\AgdaBound{n}\AgdaSymbol{)}\<%
\end{code}


\subsection{Discussion}\label{sec:q-discussion}

Everything above is checked by Agda~2.8 with standard library 2.4 under \texttt{-{}-safe}, with no postulates and no holes. The development has 215 modules and about 58{,}000 lines of code, of which about 27{,}000 are for the symplectic group, 18{,}000 for the projective chain, and 13{,}000 for the scalars. These counts exclude the shared library (\Cref{sec:compare}).

The Agda development does not formalize the two isomorphisms of \Cref{sec:q-circuits}~\cite{gross2006hudson,gurevich2007geometric}, and no matrix appears in it. The representation argument above is external mathematics, and matching the paper's gate phases remains a convention check~\cite{bian2026qupit,neuhauser2002metaplectic}. The Lean development needs neither, since it proves soundness and completeness against complex matrices directly (\Cref{sec:compare}).

We used versions of Claude for parts of this formalization: composing box relations to prove \aid{ract-sound}, defining the symplectic group, and establishing the equivalence between the code's relations and those in Figure~1 of~\quditcite{}, read modulo scalars. Claude also wrote \aid{presentation-exact} and the layer for Figure~1 with $-1$.

% \emph{Next}: a verified matrix semantics over
% $U_{p^n}(\mathbb{Z}[1/p,\omega])$, which would close the chain and
% discharge the glue facts, and submission to the Agda standard
% library.  The qubit case~\cite{selinger2015clifford} and the qutrit
% rules of~\cite{li2025qutrit}, which the odd-prime development does not
% cover, are now among the formalizations of \Cref{sec:corpus}.

\section{Formalized Completeness and Minimality Proofs}\label{sec:corpus}

The sixteen formalizations in \Cref{tab:corpus} contain about 376{,}000 lines of Agda code and use no postulates, and each theorem is as general as the result in its source paper. Moreover, all matrices are computed exactly. Their entries lie in algebraic number rings, whose elements are given by finitely many integers, so Agda checks automatically whether two matrices are equal. Over the real or complex numbers, no such automatic check exists, which limits the equality checking of QCert to discrete gate sets, as we discuss in \Cref{sec:conclusion}.



\begin{table}[tbp]
\small
\setlength{\tabcolsep}{4pt}
\renewcommand{\arraystretch}{1.3}
\begin{tabularx}{\linewidth}{@{}l Y Y r r r@{}}
\toprule
Source & Object & Result & \shortstack{Agda\\Mod} & \shortstack{Agda\\time (s)} & \shortstack{Lean\\time (s)} \\
\midrule
\multicolumn{6}{@{}l}{\emph{Groups of $m$-level matrices}}\\
\tabcite{greylyn2014} & $U_n(\Z[1/\sqrt2,i])$ & presentation, $n\le4$ & 31 & 42 & 145.34 \\
\tabcite{bian2021un} & $U_n(\Z[1/2,i])$ & presentation & 37 & 62 & 126.99 \\
\tabcite{li2021on} & $O_n(\Z[1/2])$, $L_n$ & presentations & 68 & 79 & 185.32 \\
\tabcite{fang2026hadamard} & $O_n(\Z[1/\sqrt2])$ & presentation & 44 & 44 & 231.57 \\
\tabcite{clement2019thesis} & $O_n(\Z[1/\sqrt2])$ & presentation & 40 & 64 & 51.99 \\
\midrule
\multicolumn{6}{@{}l}{\emph{Circuits}}\\
\tabcite{selinger2015clifford} & Clifford & sub-presentation & 31 & 34 & 88.24 \\
\citetalias{bian2026qupit}~[\citeyear{bian2026qupit}] & Clifford, odd prime~$p$ & presentation of a~model & 215 & 546 & 197.32 \\
\tabcite{makary2021real} & real stabilizer & sub-presentation & 30 & 35 & 79.75 \\
\tabcite{amy2018cnotdihedral} & CNOT-dihedral & sub-presentation & 30 & 24 & 101.22 \\
\tabcite{blake2026phaseaffine} & phase-affine, prime~$d$ & sub-presentations & 56 & 61 & 257.24 \\
\tabcite{bian2023cliffordt2} & two-qubit Clifford+T & sub-presentation in~$U_4$ & 62 & 118 & 356.79 \\
\tabcite{bian2023cliffordcs3} & three-qubit Clifford+CS & sub-presentation in~$U_8$ & 123 & 250 & 1330.89 \\
\tabcite{clement2026realcliffordch} & real Clifford+CH & sub-presentation & 346 & 405 & 5478.79\textsuperscript{*} \\
\tabcite{amy2024toffoli} & three-qubit Toffoli--Hadamard & presentation of $L_8$ & 201 & 102 & 134.91 \\
\tabcite{li2025qutrit} & qutrit Clifford & presentation & 89 & 282 & 179.60 \\
\midrule
\multicolumn{6}{@{}l}{\emph{Minimality}}\\
\tabcite{blake2026simpler} & six fragments & minimal, complete & 73 & 183 & 358.87 \\
\midrule
Total & & & 1{,}476 & 2{,}332 & 9{,}304.83 \\
\bottomrule
\end{tabularx}
\captionsetup{position=bottom}
\caption{Completeness and minimality results formalized in QCert. A presentation gives generators and relations for a whole group, and a sub-presentation for the subgroup generated by the gate set named in the Object column of the Circuits block. $L_n$ is the group of orthogonal matrices $N/\sqrt2^{\,j}$, where $N$ is an integer matrix and $j \ge 0$. Mod counts Agda modules. The Agda times are the previously recorded, rounded typechecking times from scratch with Agda~2.8 on an Intel Core i5-14400; the Lean times are elapsed folder-build times under the conditions of \Cref{sec:compare}. Both columns exclude the shared library and other developments on which a row depends. The Result column describes the Agda statements; \Cref{sec:compare} explains the semantic differences. \textsuperscript{*}Includes two unused legacy helpers that failed; the main theorem chain passed.}
\label{tab:corpus}
\end{table}

In the Agda development, most theorems in \Cref{tab:corpus} are proved from scratch in QCert against a matrix semantics, with three exceptions. The two rows of~\cite{bian2023cliffordt2,bian2023cliffordcs3} build on the Agda proofs that accompany those papers. These proofs relate the circuit relations to two-level relations by the Reidemeister--Schreier method~\cite{magnus2004combinatorial}. We connect them to our formalizations of the two-level matrix presentations in~\cite{greylyn2014,bian2021un}, the first two rows of the table, to obtain sub-presentations of $U_4\left(\Z\left[1/\sqrt2,i\right]\right)$ and $U_8\left(\Z\left[1/2,i\right]\right)$. The qudit Clifford theorem of~\quditcite{} is the only one whose semantics is not a matrix group. Its group presentation lifts a presentation of the symplectic group $\mathrm{Sp}(2n,\Z_p)$ with Pauli corrections and global phases, and its semantics is built in the same way (\Cref{sec:qudit}).

% Long computational derivations, such as those in Selinger's 105-page supplement and the qutrit box relations, come from scripts that follow the paper, as in \Cref{sec:method}, or search for the steps. In the real stabilizer formalization, half the lines are generated in this way, and a fifth formalize the paper's mathematics.

\section{Corrections to the Source Papers}\label{sec:findings}

The formalizations in \Cref{tab:corpus} identify two kinds of issues in the source papers. Most are typos, and some of them appear in statements that the proofs rely on. These are easy to fix. Two others are harder to repair. One is a gap in a proof, and the other is a relation missing from a theorem. We describe both below.

The gap is in the proof of Theorem~4.14 of Hadamard-$\Pi$~\cite{fang2026hadamard}, which states that the relations in its Figure~6, over the one- and two-level matrices $(-1)_{[a]}$, $\mathrm{X}_{[a,b]}$, and $\mathrm{H}_{[a,b]}$, present the group $O_n\left(\Z\left[1/\sqrt2\right]\right)$. One case of the proof relies on Lemmas~A.16 and~A.17, which have counterexamples in dimension~$6$, and the case analysis built on them misses a matrix in dimension~$9$. The authors have confirmed this gap (personal communication). 
Their main result is however still correct: Cl\'ement's unpublished notes~\cite{clement2019thesis} present the same group with different relations, and his case analysis covers the missing case, using which the formalization of Theorem~4.14 is completed. We also prove that the two sets of relations are equivalent in every dimension and formalize Cl\'ement's proof~\cite{qcert2026clement}.

The relation is missing from Theorem~5.2 of~\cite{amy2024toffoli}. It states that the relations in its Figures~4 and~7, over the Toffoli and $\mathrm{K} = \mathrm{H}\otimes\mathrm{H}$ gates, present the group $O_8\left(\Z\left[1/2\right]\right)$. This is the first of two steps towards presenting the scaled dyadic orthogonal group $L_8$ by three-qubit Toffoli--Hadamard circuits. The relation $r_{\mathrm{K}}$ in its Section~5.5, which splits $\mathrm{K}$ on two qubits into two copies controlled by the third qubit (\Cref{sec:previous-completeness}), is in neither figure. Nor can it be derived from them. Hence, Theorem~5.2 fails as stated, and we have informed the authors. Our formalization~\cite{qcert2026toffoli} shows that $r_{\mathrm{K}}$ is independent by giving an interpretation of circuits that every other relation preserves but $r_{\mathrm{K}}$ does not, and proves the presentations of $O_8\left(\Z\left[1/2\right]\right)$ and $L_8$ with $r_{\mathrm{K}}$ restored.


\section{Discussion and Future Work}\label{sec:conclusion}

We formalize the results of sixteen papers in QCert, using consistent circuit syntax, a unified definition of group presentations, and shared libraries for rings and their matrix groups. Since all formalizations use these definitions, presentations of the same group from different papers can be compared and proved equivalent. The same libraries can be reused for new algorithms, and the induction that turns an exact synthesis algorithm into a completeness proof carries over with little change. The main part specific to each algorithm is the case analysis over matrices and generators. As the most error-prone part of the proof, it is where machine checking in QCert matters most.

At the moment, not every completeness result fits into QCert, as its scope is limited in two ways. First, it supports only discrete gate sets, so it does not cover completeness results for circuits with continuous parameters, such as rotations by arbitrary angles~\cite{clement2023complete,clement2024extensions,clement2024minimal}. Second, control is not a circuit constructor in QCert, so it does not cover results that treat control as a primitive operation on circuits~\cite{delorme2026control,schober2026control,blake2026qudit}.

We plan to extend the infrastructure of QCert to more families of quantum circuits, including the two above. Completeness is still open for many discrete gate sets in higher-dimensional quantum systems, and their case analyses are often too large to check by hand. In particular, we plan to study the completeness of Clifford circuits in all dimensions with the help of QCert. Beyond completeness, the exact synthesis libraries can be reused for T-count optimization of multi-qubit Clifford+T circuits. 

\bibliographystyle{ACM-Reference-Format}
\bibliography{refs}

\clearpage
\appendix
\iffalse
\section{A New Result: Toffoli–Hadamard Completeness for Any Number of Qubits}
\label{sec:toffoli-hadamard}

The Toffoli--Hadamard gate set $\{\mathrm{X},\mathrm{CX},\mathrm{CCX},\mathrm{H}\}$ adds the Hadamard gate to the classical reversible gates $\mathrm{X}$, $\mathrm{CX}$, and $\mathrm{CCX}$. It is computationally universal, meaning that it can efficiently simulate arbitrary quantum computations to arbitrary accuracy~\cite{shi2003both,aharonov2003simple}. This makes it useful in applications such as Grover-based AES key search~\cite{grassl2016applying} and fault-tolerant computation using transversal gates and error correction~\cite{paetznick2013universal}. Toffoli--Hadamard circuits also have a simple arithmetic structure, which makes them useful for exact synthesis and verification. Their matrices are orthogonal and have a common denominator that is a power of $\sqrt2$, a property exploited by exact synthesis algorithms~\cite{amy2020numbertheoretic,amy2023improved}. A complete set of rewrite rules for these circuits would support symbolic proofs of circuit equivalence and provide a basis for verifying circuit optimizations. In this section, we seek a complete theory stated directly as circuit equations, on any number of qubits and without ancillas.
%Aharonov's proof starts from a known universal gate set: Hadamard and controlled-$\mathrm{S}$, where $\mathrm{S}=\operatorname{diag}(1,i)$. The simulation uses one additional qubit to encode complex amplitudes as real amplitudes. Hadamard gates act unchanged, and each encoded controlled-$\mathrm{S}$ gate is implemented by two Toffoli gates and two Hadamard gates~\cite{aharonov2003simple}.

\subsection{Preliminaries}
\label{sec:preliminaries}

Besides Toffoli--Hadamard circuits, we consider the Toffoli--K gate set $\{\mathrm{X}, \mathrm{CX}, \mathrm{CCX}, \mathrm{K}\}$, where $\mathrm{K} = \mathrm{H}\otimes\mathrm{H}$~\cite{amy2020numbertheoretic,li2021on,amy2024toffoli}. Recall that
\[
\mathrm{X} = \begin{pmatrix} 0 & 1 \\ 1 & 0 \end{pmatrix}, \qquad
\mathrm{H} = \frac{1}{\sqrt2}\begin{pmatrix} 1 & 1 \\ 1 & -1 \end{pmatrix}, \qquad
\mathrm{K} = \frac{1}{2}\begin{pmatrix} 1 & 1 & 1 & 1 \\ 1 & -1 & 1 & -1 \\ 1 & 1 & -1 & -1 \\ 1 & -1 & -1 & 1 \end{pmatrix}.
\]
We number the qubits from $0$, with qubit~$0$ the most significant, so the basis state $\lvert x_0x_1x_2\rangle$ of three qubits is the basis vector $e_{4x_0+2x_1+x_2}$. We write $\mathrm{X}_j$ and $\mathrm{H}_j$ for the $\mathrm{X}$ and $\mathrm{H}$ gates on qubit~$j$, and $\mathrm{K}_{j,k}$ for $\mathrm{H}_j\mathrm{H}_k$.

The interpretations of these circuits are orthogonal matrices. We write $O_n\left(\Z\left[1/2\right]\right)$ for the group of $n\times n$ orthogonal matrices with entries in the ring of dyadic rationals $\Z\left[1/2\right]$. We write $L_n$ for the group of $n\times n$ orthogonal matrices of the form $N/\sqrt2^{\,j}$, where $N$ is an integer matrix and $j\ge0$. The first group is a subgroup of the second, of index~$2$ when $n$ is even~\cite{li2021on}. On $n$ qubits, the interpretations of Toffoli--K circuits lie in $O_{2^n}\left(\Z\left[1/2\right]\right)$, and those of Toffoli--Hadamard circuits lie in $L_{2^n}$. With ancillas, every matrix in these groups is the interpretation of some circuit~\cite{amy2020numbertheoretic}. Without ancillas, as in this section, the circuits may generate only a subgroup. Hence, we state completeness on $n$ qubits as a sub-presentation (see \Cref{tab:corpus}).

\citet{li2021on} gave presentations of both groups in every dimension~$n$. Their generators are multi-level matrices (\Cref{sec:background}) rather than gates. The generators of $O_n\left(\Z\left[1/2\right]\right)$ are the one-, two-, and four-level matrices
\[
(-1)_{[a]}, \qquad \mathrm{X}_{[a,b]}, \qquad \mathrm{K}_{[a,b,c,d]}, \qquad 0\le a<b<c<d<n,
\]
and a finite list of relation schemas, stated independently of~$n$, presents this group~\cite[Theorem~4.10]{li2021on}. For even~$n$, adding the generator $\mathrm{I}\otimes\mathrm{H} = \mathrm{diag}(\mathrm{H},\dots,\mathrm{H})$ and four relations gives a presentation of $L_n$~\cite[Theorem~5.7]{li2021on}. A few generators are gates. On three qubits, for example, the Toffoli gate controlled by qubits~$0$ and~$1$ is $\mathrm{X}_{[6,7]}$, and $\mathrm{I}\otimes\mathrm{H}$ is $\mathrm{H}_2$. Most generators, however, are not gates of either gate set. For instance, $\mathrm{K}_{[4,5,6,7]}$ applies $\mathrm{K}_{1,2}$ only when qubit~$0$ is in the state $\lvert 1\rangle$. So these presentations are not yet equational theories for circuits.
\iffalse
\subsection{Previous Completeness Proofs}
\label{sec:previous-completeness}

In 2021, \citet{li2021on} found a complete set of rewrite rules for the group $O_n\left(\Z\left[1/2\right]\right)$, whose generators are one-, two-, and four-level matrices. These generators are matrices rather than gates, so the rules did not yet apply to circuits. In 2023, \citet{vilmart2023completeness} gave a complete set of rewrite rules for Toffoli--Hadamard computations on any number of qubits, in the sum-over-paths formalism. Its terms, however, represent general linear maps rather than circuits. Completeness for such a class of linear maps does not in general yield complete rewrite rules for circuits, because a derivation between two circuits may pass through terms that are not circuits. In 2024, \citet{amy2024toffoli} used the previously known fact that on three qubits, the gates $\mathrm{X}$, $\mathrm{CX}$, $\mathrm{CCX}$, and $\mathrm{H}\otimes\mathrm{H}$ generate a finite group, the automorphism group of the $E_8$ lattice. They first turned its well-known Coxeter presentation into rewrite rules for these circuits. They then used the rules of~\citet{li2021on} for $n=8$ to add the gates $\mathrm{CCZ}$ and $\mathrm{H}$. This gave a complete set of rewrite rules for three-qubit Toffoli--Hadamard circuits, once a missing relation is restored (\Cref{sec:findings}). Their method stops at three qubits. On four or more qubits, the same gates generate an infinite group, so there is no finite group to start from. 

\Cref{fig:amy-completeness} details the proof in three steps. \emph{Step~1} (\Cref{fig:amy-toffoli-k}): translate Toffoli--$\mathrm K$ circuits $U_j$ with equal matrices into reflection sequences $F_j=E_0(U_j)$. The known complete rules $\mathcal R_{E_8}$ connect $F_1$ and $F_2$. Prove that each reflection rule follows from circuit rules $\mathcal R_0$, and that translating back recovers the original circuits. Thus $D_0$ turns the reflection proof into a circuit proof, establishing completeness of $\mathcal R_0$. \emph{Step~2} (\Cref{fig:amy-dyadic-transfer}): add $\mathrm{CCZ}$ and a helper that applies $\mathrm K$ to qubits 1 and 2 when qubit 0 is $0$. These circuits $B_j$ represent $O_8(\D)$, where $\D=\Z[1/2]$. Translate them into matrix-operation sequences $M_j=E_1(B_j)$. Derive Li et al.'s complete rules $\mathcal R_{\mathrm{Li}}$~\cite{li2021on} from $\mathcal R_0$ and additional circuit rules, including the restored relation
\[
r_{\mathrm K}:\quad \mathrm{K}_{1,2} = \mathrm{K}_{[4,5,6,7]}\,\mathrm{X}_0\,\mathrm{K}_{[4,5,6,7]}\,\mathrm{X}_0.
\]
Here $\mathrm{K}_{[4,5,6,7]}$ applies $\mathrm{K}_{1,2}$ when qubit~$0$ is $\lvert1\rangle$, and the helper $\mathrm{X}_0\,\mathrm{K}_{[4,5,6,7]}\,\mathrm{X}_0$ applies it when qubit~$0$ is $\lvert0\rangle$, so together they apply $\mathrm{K}_{1,2}$. Translating proofs back by $D_1$ gives completeness of the resulting rules $\mathcal R_1$~\cite[Theorem~5.2, with $r_{\mathrm K}$ restored]{amy2024toffoli}. \emph{Step~3} (\Cref{fig:amy-hadamard-extension}): extend $\mathcal R_1$ to $\mathcal R_2$ with Hadamard rules. These move Hadamards to the end, changing intervening gates, and cancel pairs, giving $C_j\rewrites{\mathcal R_2}B_jH_2^{\epsilon}$, where $H_2$ acts on qubit 2 and $\epsilon\in\{0,1\}$. Equal matrices force the same $\epsilon$: $B_j$ has rational entries, whereas $B_jH_2$ has an irrational entry. Cancel the common final factor, connect $B_1$ and $B_2$ using $\mathcal R_1$, and restore the factor. Hence $\mathcal R_2$ is complete~\cite[Theorem~6.3]{amy2024toffoli}.

\begin{figure}[!thb]
\centering
\begin{subfigure}[t]{0.48\linewidth}
  \centering
  \scalebox{.82}{\tikzfig{figures/amy-toffoli-k}}
  \caption{$U_j$: Toffoli--$\mathrm K$ circuits; $F_j$: $E_8$ reflection sequences.}
  \label{fig:amy-toffoli-k}
\end{subfigure}\hfill
\begin{subfigure}[t]{0.48\linewidth}
  \centering
  \scalebox{.82}{\tikzfig{figures/amy-dyadic-transfer}}
  \caption{$B_j$: circuits for $O_8(\D)$; $M_j$: matrix-operation sequences.}
  \label{fig:amy-dyadic-transfer}
\end{subfigure}

\medskip
\begin{subfigure}[t]{0.60\linewidth}
  \centering
  \scalebox{.82}{\tikzfig{figures/amy-hadamard-extension}}
  \caption{$C_j$: Toffoli--Hadamard circuits; $B_jH_2^{\epsilon}$: at most one final Hadamard.}
  \label{fig:amy-hadamard-extension}
\end{subfigure}
\caption{Three stages of completeness for three-qubit Toffoli--Hadamard circuits~\cite{amy2024toffoli}, with $r_{\mathrm K}$ restored (\Cref{sec:findings}). Vertices label circuits or operator sequences. Dashed $\equiv_u$ lines assert equal matrices; $\mathcal R^{*}$ arrows denote zero or more rewrites in either direction. $E_i$ translates descriptions, and $D_i$ translates proofs back into circuit rewrites.}
\Description{Three proof diagrams. The first translates paired-Hadamard circuits into E8 reflections and back. The second translates rational-entry circuits into Li et al.'s matrix operations and back. The third rewrites full circuits into a rational-entry part followed by at most one Hadamard, then uses the second stage on their rational-entry parts.}
\label{fig:amy-completeness}
\end{figure}

Recently, \citet{clement2026realcliffordch} found a complete set of rewrite rules on any number of qubits for a closely related family, the real Clifford+CH circuits over the gates $\mathrm{H}$, $\mathrm{Z}$, $\mathrm{CZ}$, and the controlled Hadamard gate $\mathrm{CH}$. His rules consist of eighteen equations on at most three qubits and one equation schema, which gives an equation on $n$ qubits for each $n\ge5$. The rules we find for Toffoli--Hadamard circuits in \Cref{sec:arbitrary-qubit-completeness} have the same shape, a few equations on a small number of qubits together with one schema over the number of qubits. Our proof also follows his. He translates circuits back and forth to a presentation by one- and two-level matrices, as in \Cref{fig:translation}, and we do the same with the presentations of~\citet{li2021on}.
\fi
\subsection{Completeness for Any Number of Qubits}
\label{sec:arbitrary-qubit-completeness}
\fi

\end{document}
